% =============================================================================
%  Simulator Methodology — Thesis Section
%  To include in a chapter: % =============================================================================
%  Simulator Methodology — Thesis Section
%  To include in a chapter: % =============================================================================
%  Simulator Methodology — Thesis Section
%  To include in a chapter: % =============================================================================
%  Simulator Methodology — Thesis Section
%  To include in a chapter: \input{simulator_methodology}
% =============================================================================

\section{Simulator Implementation}
\label{sec:simulator_implementation}

This section describes how the trajectory simulation tool is structured and how
the theoretical principles introduced in the previous chapter are brought
together into a working ascent simulator.  The discussion covers the overall
software architecture, the reference frames adopted, the equations of motion
actually integrated, the environmental and force models, the guidance laws
available, and the optimisation strategy used to find the best initial
conditions.  Vehicle parameters representative of a Falcon~9-class launcher
close the section.

% -----------------------------------------------------------------------------
\subsection{Simulation Architecture and Workflow}
\label{ssec:sim_architecture}

The simulator is organised around three successive mission phases, each
resolved by a separate call to the numerical integrator:

\begin{enumerate}
  \item \textbf{Stage~1 powered ascent.}  From lift-off until first-stage
        propellant depletion (Main Engine Cut-Off, MECO).  During this phase
        the vehicle follows an initial vertical rise, a pitchover kick
        manoeuvre, and atmospheric flight with aerodynamic drag.

  \item \textbf{Stage~2 powered ascent with guidance.}  After stage separation
        and a brief coast interval, the second-stage engine ignites.  An active
        guidance law (selected from five available modes) shapes the trajectory
        until the event-triggered Second Engine Cut-Off (SECO), which occurs
        when the instantaneous apogee of the coasting trajectory matches the
        target orbital altitude.

  \item \textbf{Ballistic coast and circularisation.}  The vehicle coasts
        unpowered to apogee, where an impulsive circularisation burn is applied
        to achieve the target circular orbit.  A final short coast in stable
        orbit concludes the simulation.
\end{enumerate}

\subsubsection{State Vector}

The base state vector integrated by the solver is
\begin{equation}
  \mathbf{x} = \bigl[\,s,\; r,\; v,\; \gamma,\; m\,\bigr]^T,
  \label{eq:state_base}
\end{equation}
where $s$ is the downrange distance along the Earth surface, $r$ is the
geocentric radius, $v$ is the velocity magnitude in the rotating frame,
$\gamma$ is the flight-path angle, and $m$ is the vehicle mass.

When Earth-rotation effects are enabled, the state is extended to include
geocentric latitude~$\phi$:
\begin{equation}
  \mathbf{x} = \bigl[\,s,\; r,\; v,\; \gamma,\; m,\; \phi\,\bigr]^T.
  \label{eq:state_rot}
\end{equation}

If heading tracking is additionally activated, the local heading (azimuth)
$\psi$ becomes the seventh state component:
\begin{equation}
  \mathbf{x} = \bigl[\,s,\; r,\; v,\; \gamma,\; m,\; \phi,\; \psi\,\bigr]^T.
  \label{eq:state_heading}
\end{equation}

\subsubsection{Numerical Integration}

Each mission phase is propagated with an explicit Runge--Kutta~4(5) scheme
(Dormand--Prince variant) provided by \texttt{scipy.integrate.solve\_ivp}.  The
solver uses adaptive step-size control with an absolute tolerance of
$10^{-8}$, a maximum integration step of $1$~s, and a dense-output time
resolution of $0.01$~s.

Phase transitions are handled through \emph{event functions}: scalar functions
of the state that the integrator monitors for zero-crossings.  When a crossing
is detected the integrator performs root-finding to locate the exact transition
time and terminates the current phase.  Events used include propellant
depletion (MECO), stage separation timing, atmosphere exit, target-apogee
matching (SECO), and ground-collision safeguards.

\subsubsection{Architecture Overview}

Figure~\ref{fig:sim_architecture} shows the high-level data flow.  An outer
optimisation loop varies the initial kick angle supplied to the simulation
engine.  Inside each simulation call, the equations of motion are evaluated at
every integrator step using the current environmental models (gravity,
atmosphere) and the active guidance law to compute the steering command.  Event
functions run concurrently and trigger phase switches when their conditions are
met.  After the full trajectory is obtained, the optimiser evaluates the
objective function and proceeds to the next candidate.

\begin{figure}[htbp]
  \centering
  \begin{tikzpicture}[
      >=latex,
      block/.style  = {draw, rounded corners, minimum height=1.0cm,
                        minimum width=2.6cm, align=center, font=\small},
      smallblock/.style = {draw, rounded corners, minimum height=0.8cm,
                        minimum width=2.2cm, align=center, font=\footnotesize},
      arrow/.style  = {->, thick},
      darrow/.style = {<->, thick},
    ]

    % Optimisation loop
    \node[block, fill=blue!12] (opt) {Optimisation\\Loop};

    % Simulation engine
    \node[block, fill=green!10, right=2.2cm of opt] (sim) {Simulation\\Engine};

    % Inner modules (below sim)
    \node[smallblock, fill=orange!10, below left=1.4cm and -0.3cm of sim]
          (guid) {Guidance\\Module};
    \node[smallblock, fill=orange!10, below=1.4cm of sim]
          (eom) {Equations\\of Motion};
    \node[smallblock, fill=orange!10, below right=1.4cm and -0.3cm of sim]
          (env) {Environment\\Models};

    % Event detection
    \node[block, fill=red!10, right=2.2cm of sim] (evt) {Event\\Detection};

    % Outputs
    \node[block, fill=gray!12, right=2.2cm of evt] (out) {Trajectory\\Outputs};

    % Arrows
    \draw[arrow] (opt) -- node[above, font=\footnotesize]{$\alpha_{\text{kick}}$} (sim);
    \draw[arrow] (sim) -- (evt);
    \draw[arrow] (evt) -- (out);
    \draw[arrow] (out.south) -- ++(0,-1.2) -| node[near start, below, font=\footnotesize]{$J(\alpha)$} (opt.south);

    % Inner connections
    \draw[arrow] (sim) -- (guid);
    \draw[arrow] (sim) -- (eom);
    \draw[arrow] (sim) -- (env);
    \draw[arrow] (guid) -- (eom);
    \draw[arrow] (env) -- (eom);

  \end{tikzpicture}
  \caption{High-level architecture of the trajectory simulator.  The
           optimisation loop feeds candidate kick angles into the simulation
           engine, which dispatches environment evaluations and guidance
           commands to the equations of motion at each integration step.  Event
           detection triggers phase transitions and engine cutoff.  The
           resulting trajectory is scored and fed back to the optimiser.}
  \label{fig:sim_architecture}
\end{figure}

% -----------------------------------------------------------------------------
\subsection{Reference Frames}
\label{ssec:reference_frames}

The simulator employs a mixed-frame strategy that keeps guidance and
atmosphere coupling straightforward while preserving correct orbital-mechanics
metrics.

\begin{itemize}
  \item \textbf{Rotating frame (ECEF-like).}  All ascent dynamics quantities
        --- velocity magnitude~$v$, flight-path angle~$\gamma$, and heading
        decomposition --- are expressed in a frame that co-rotates with the
        Earth.  This simplifies the treatment of aerodynamic drag, which
        depends on the velocity relative to the atmosphere.

  \item \textbf{Inertial frame (ECI).}  Orbital-energy and orbital-element
        calculations (semi-major axis, eccentricity, apogee radius) are
        performed in an Earth-Centred Inertial frame.  This is essential for
        correct event detection (e.g.\ target-apogee matching) and end-of-run
        orbit characterisation.

  \item \textbf{Local East--North--Up (ENU).}  Pseudo-force accelerations
        (Coriolis and centrifugal) are first evaluated in the local ENU basis at
        the vehicle's current latitude, then projected onto the along-heading
        and radial directions used by the equations of motion.
\end{itemize}

\subsubsection{Frame Transition During Flight}

During powered ascent and early coasting the integration proceeds in the
rotating frame with optional pseudo-force corrections.  Once the vehicle
reaches the final ballistic coast phase (post-SECO), the propagation switches
to pure inertial dynamics: pseudo-force terms are deactivated to avoid
introducing spurious accelerations when the trajectory is governed solely by
two-body gravity.

\subsubsection{ECEF-to-ECI Velocity Conversion}

At any instant the local rotating-frame velocity components are converted to
their inertial counterparts by adding the Earth's rotational eastward speed.
The horizontal velocity magnitude is $v_h = v\cos\gamma$, which decomposes as
\begin{equation}
  v_N = v_h\cos\psi, \qquad
  v_E^{\text{ECEF}} = v_h\sin\psi.
  \label{eq:vel_ecef}
\end{equation}

Adding Earth rotation yields
\begin{equation}
  v_E^{\text{ECI}} = v_E^{\text{ECEF}} + \Omega_E\, r\cos\phi,
  \label{eq:ve_eci}
\end{equation}
from which the inertial speed and flight-path angle follow:
\begin{align}
  v_{\text{ECI}} &= \sqrt{\bigl(v_E^{\text{ECI}}\bigr)^2 + v_N^2
                    + \bigl(v\sin\gamma\bigr)^2}, \label{eq:v_eci} \\[4pt]
  \gamma_{\text{ECI}} &= \operatorname{atan2}\!\bigl(v\sin\gamma,\;
                    \sqrt{\bigl(v_E^{\text{ECI}}\bigr)^2 + v_N^2}\;\bigr).
                    \label{eq:gamma_eci}
\end{align}

These inertial quantities are used whenever orbital elements are computed.

% -----------------------------------------------------------------------------
\subsection{Launch Azimuth and Heading}
\label{ssec:azimuth_heading}

\subsubsection{Launch Azimuth Determination}
\label{sssec:launch_azimuth}

The launch azimuth defines the initial heading of the ascent trajectory in the
horizontal plane and directly controls the achievable orbital inclination.
Two modes are available.

\paragraph{Geometric (inertial) azimuth.}
Given a target inclination~$i$ and launch latitude~$\phi$, the classical
spherical-geometry relation gives the inertial azimuth measured from true
north:
\begin{equation}
  \sin\beta_{\text{geom}} = \frac{\cos i}{\cos\phi}.
  \label{eq:az_geom}
\end{equation}

\paragraph{Corrected (rotating-frame) azimuth.}
Because the rotating-frame velocity seen by the guidance system differs from
the inertial velocity by the eastward surface speed, the launch azimuth is
corrected:
\begin{equation}
  v_{\text{rot}} = \Omega_E\, R_E\cos\phi,
  \label{eq:vrot}
\end{equation}
\begin{align}
  v_N &= v_{\text{orb}}\cos\beta_{\text{geom}}, \label{eq:vn_corr} \\
  v_E^{\text{ECEF}} &= v_{\text{orb}}\sin\beta_{\text{geom}} - v_{\text{rot}},
  \label{eq:ve_corr}
\end{align}
where $v_{\text{orb}} = \sqrt{\mu/r_{\text{target}}}$ is the circular speed at
the target altitude.  The corrected azimuth is then
\begin{equation}
  \beta_{\text{corr}} = \operatorname{atan2}\!\bigl(v_E^{\text{ECEF}},\, v_N\bigr).
  \label{eq:az_corr}
\end{equation}

The simulator defaults to the corrected mode; a configuration flag may select
the geometric mode instead.

\subsubsection{Heading Propagation}
\label{sssec:heading_propagation}

When heading tracking is enabled, the local azimuth $\psi$ is propagated as a
state variable.  Its rate is computed from spherical-geometry consistency with
the great-circle ground track:
\begin{equation}
  \dot\psi = \tan\phi\;\tan\psi\;\dot\phi,
  \label{eq:psi_dot}
\end{equation}
where $\dot\phi$ is obtained from the latitude kinematic relation detailed in
Section~\ref{ssec:eom}.  An optional cross-heading pseudo-force contribution
may be added:
\begin{equation}
  \dot\psi \mathrel{+}= \frac{a_{\perp}}{v\cos\gamma},
  \label{eq:psi_pseudo}
\end{equation}
with $a_{\perp}$ the component of the Coriolis and centrifugal acceleration
perpendicular to the heading direction (see Section~\ref{sssec:pseudo_forces}).

\subsubsection{Achieved Inclination}
\label{sssec:achieved_inclination}

At the end of the simulation the achieved orbital inclination is evaluated by
converting the final rotating-frame velocity to inertial components using
Eqs.~\eqref{eq:vel_ecef}--\eqref{eq:ve_eci} and computing the inertial
velocity azimuth:
\begin{equation}
  \beta_{\text{ECI}} = \operatorname{atan2}\!\bigl(v_E^{\text{ECI}},\, v_N\bigr).
  \label{eq:beta_eci}
\end{equation}

The inclination follows from the classical relation:
\begin{equation}
  \cos i_{\text{achieved}} = \cos\phi\;\sin\beta_{\text{ECI}}.
  \label{eq:inc_achieved}
\end{equation}

The inclination drift $\Delta i = i_{\text{achieved}} - i_{\text{target}}$ is
a valuable metric for assessing the fidelity of the rotating-frame treatment
and the effectiveness of the heading-tracking augmentation.

% -----------------------------------------------------------------------------
\subsection{Equations of Motion}
\label{ssec:eom}

\subsubsection{Base Dynamics (Non-Rotating Form)}

The core differential system, expressed in spherical coordinates before
optional pseudo-force augmentation, is
\begin{align}
  \dot s      &= \frac{R_E}{r}\,v\cos\gamma,
               \label{eq:sdot}\\[3pt]
  \dot r      &= v\sin\gamma,
               \label{eq:rdot}\\[3pt]
  \dot v      &= \frac{F_T}{m}\cos\alpha
                 - \frac{F_D}{m}
                 - g(r)\sin\gamma,
               \label{eq:vdot}\\[3pt]
  \dot\gamma  &= \frac{1}{v}\!\left(
                   \frac{F_T}{m}\sin\alpha
                   + \frac{F_L}{m}
                   - \Bigl[g(r) - \frac{v^2}{r}\Bigr]\cos\gamma
                 \right),
               \label{eq:gammadot}\\[3pt]
  \dot m      &= -\frac{F_T}{I_{sp}\,g_0},
               \label{eq:mdot}
\end{align}
where $\alpha$ is the angle of attack (steering command from the guidance law),
$F_T$ and $I_{sp}$ are the current thrust and specific impulse, $F_D$ is
aerodynamic drag, $F_L$ is lift (set to zero in the present implementation),
and $g(r) = \mu/r^2$ is the gravitational acceleration.

The term $v^2/r$ in Eq.~\eqref{eq:gammadot} represents the centripetal
acceleration arising from the curved trajectory; when $v^2/r$ exceeds
$g\cos\gamma$ the net effect tends to raise the flight-path angle.

Two numerical safeguards are applied:
\begin{itemize}
  \item \emph{Small-velocity guard:} when $v < 10^{-6}$\,m/s the angular rate
        $\dot\gamma$ is forced to zero to prevent a singularity.
  \item \emph{Vertical-hold constraint:} during the initial vertical rise
        (before the pitchover kick begins) $\dot\gamma$ is explicitly set to
        zero regardless of geometric terms.
\end{itemize}

\subsubsection{Rotating-Frame Augmentation}

When Earth-rotation effects are enabled, Coriolis and centrifugal
accelerations computed in the local ENU frame (see
Section~\ref{sssec:pseudo_forces}) are projected onto the along-heading and
radial directions:
\begin{equation}
  a_{h,\parallel} = a_E\sin\psi + a_N\cos\psi,
  \label{eq:a_parallel}
\end{equation}
where $a_E$ and $a_N$ collect both Coriolis and centrifugal components.
The additive corrections to the velocity and flight-path-angle rates are
\begin{align}
  \Delta\dot v     &= a_{h,\parallel}\cos\gamma + a_{\text{up}}\sin\gamma,
                    \label{eq:dv_pseudo}\\[3pt]
  \Delta\dot\gamma &= \frac{-a_{h,\parallel}\sin\gamma
                      + a_{\text{up}}\cos\gamma}{v},
                    \label{eq:dgamma_pseudo}
\end{align}
so that the integrated rates become
$\dot v_{\text{tot}} = \dot v + \Delta\dot v$ and
$\dot\gamma_{\text{tot}} = \dot\gamma + \Delta\dot\gamma$.

\subsubsection{Latitude Propagation}

Rather than integrating a differential latitude rate directly, the simulator
reconstructs latitude from the downrange angle using great-circle geometry,
which avoids numerical drift:
\begin{equation}
  \sigma = \frac{s}{R_E}, \qquad
  \sin\phi = \sin\phi_0\cos\sigma + \cos\phi_0\sin\sigma\cos\beta_0,
  \label{eq:lat_recon}
\end{equation}
where $\phi_0$ is the launch latitude and $\beta_0$ is the inertial launch
azimuth.  The latitude rate required by the heading equation
\eqref{eq:psi_dot} is obtained by differentiating analytically:
\begin{equation}
  \dot\phi = \frac{-\sin\phi_0\sin\sigma
    + \cos\phi_0\cos\sigma\cos\beta_0}{\cos\phi}\;\frac{\dot s}{R_E}.
  \label{eq:phidot}
\end{equation}

% -----------------------------------------------------------------------------
\subsection{Environmental Models}
\label{ssec:env_models}

\subsubsection{Gravitational Model}
\label{sssec:gravity}

Gravity is modelled as a central inverse-square field:
\begin{equation}
  g(r) = \frac{\mu}{r^2},
  \label{eq:gravity}
\end{equation}
with $\mu = 3.986\times10^{14}$\,m$^3$\,s$^{-2}$.  The Earth is assumed
spherical; higher-order harmonics such as $J_2$ are not included.

\subsubsection{Atmospheric Model}
\label{sssec:atmosphere}

Atmospheric density follows an exponential profile:
\begin{equation}
  \rho(h) = \rho_0\,\exp\!\left(-\frac{h}{H}\right),
  \qquad h = r - R_E,
  \label{eq:atmosphere}
\end{equation}
with sea-level density $\rho_0 = 1.225$\,kg\,m$^{-3}$ and scale height
$H = 8500$\,m.  The dynamic pressure is $q = \tfrac{1}{2}\rho\,v^2$.

Two criteria are available for declaring atmosphere exit and
activating exo-atmospheric guidance:
\begin{itemize}
  \item \emph{Altitude threshold:} $h > 65$\,km.
  \item \emph{Dynamic-pressure threshold:} $q < 1000$\,Pa (configurable).
\end{itemize}
The selected criterion is a simulation parameter.

% -----------------------------------------------------------------------------
\subsection{External Forces}
\label{ssec:forces}

\subsubsection{Aerodynamic Forces}
\label{sssec:aero}

The aerodynamic drag force opposing the velocity vector is
\begin{equation}
  F_D = q\,C_D\,A = \tfrac{1}{2}\,\rho\,v^2\,C_D\,A,
  \label{eq:drag}
\end{equation}
with constant drag coefficient $C_D = 0.3$ and reference area $A = 10.52$\,m$^2$.
A lift model is available in the formulation (Eq.~\ref{eq:gammadot}), but in
all analyses presented here lift is neglected ($F_L = 0$).

\subsubsection{Propulsive Forces}
\label{sssec:propulsion}

Thrust magnitude $F_T$ and specific impulse $I_{sp}$ are stage-dependent
constants that change at discrete events (MECO, stage separation, second-stage
ignition, SECO).  Thrust is directed at an angle of attack $\alpha$ relative
to the velocity vector; $\alpha$ is the primary control variable and is
determined at each integration step by the active guidance law
(Section~\ref{ssec:guidance}).

The resulting vehicle mass depletion obeys
\begin{equation}
  \dot m = -\frac{F_T}{I_{sp}\,g_0},
  \label{eq:mdot_prop}
\end{equation}
where $g_0 = 9.80665$\,m\,s$^{-2}$ is the standard gravitational acceleration.

\subsubsection{Pseudo-Forces (Earth Rotation)}
\label{sssec:pseudo_forces}

When the integration is performed in the rotating frame with Earth-rotation
corrections enabled, inertial pseudo-accelerations are applied.

\paragraph{Coriolis acceleration.}
In the local ENU basis at latitude $\phi$:
\begin{align}
  a_E^{\text{cor}} &= -2\Omega_E\cos\phi\;v_{\text{up}}
                      + 2\Omega_E\sin\phi\;v_N,
                    \label{eq:cor_e} \\
  a_N^{\text{cor}} &= -2\Omega_E\sin\phi\;v_E,
                    \label{eq:cor_n} \\
  a_{\text{up}}^{\text{cor}} &= 2\Omega_E\cos\phi\;v_E.
                    \label{eq:cor_u}
\end{align}

\paragraph{Centrifugal acceleration.}
\begin{align}
  a_E^{\text{cen}} &= 0,
                    \label{eq:cen_e} \\
  a_N^{\text{cen}} &= -\Omega_E^2\,r\,\sin\phi\cos\phi,
                    \label{eq:cen_n} \\
  a_{\text{up}}^{\text{cen}} &= \Omega_E^2\,r\,\cos^2\!\phi.
                    \label{eq:cen_u}
\end{align}

\paragraph{Projection onto flight-dynamics directions.}
The total pseudo-acceleration in each ENU direction is summed ($a_E = a_E^{\text{cor}} + a_E^{\text{cen}}$, etc.) and then projected onto the heading direction and the vertical:
\begin{align}
  a_{h,\parallel} &= a_E\sin\psi + a_N\cos\psi
    & &\text{(along-heading)}, \label{eq:a_par} \\
  a_{h,\perp}     &= a_E\cos\psi - a_N\sin\psi
    & &\text{(cross-heading)}, \label{eq:a_perp}
\end{align}
with $a_{\text{up}}$ retained for the radial direction.

The along-heading and up components enter the velocity and flight-path-angle
corrections (Eqs.~\ref{eq:dv_pseudo}--\ref{eq:dgamma_pseudo}).  If cross-heading
pseudo-force tracking is enabled, $a_{h,\perp}$ contributes to the heading
rate via Eq.~\eqref{eq:psi_pseudo}.

% -----------------------------------------------------------------------------
\subsection{Performance Metrics}
\label{ssec:performance}

\subsubsection{Orbital Elements}
\label{sssec:orbital_elements}

Once the velocity is converted to the inertial frame
(Eqs.~\ref{eq:v_eci}--\ref{eq:gamma_eci}), the standard two-body orbital
elements are computed:
\begin{align}
  a &= \frac{\mu\,r}{2\mu - r\,v_{\text{ECI}}^2},
    \label{eq:sma}\\[3pt]
  e &= \sqrt{1 - \frac{\bigl(r\,v_{\text{ECI}}\cos\gamma_{\text{ECI}}\bigr)^2}
       {\mu\,a}},
    \label{eq:ecc}\\[3pt]
  r_{\text{apo}} &= a\,(1+e), \qquad
  r_{\text{peri}} = a\,(1-e).
  \label{eq:apsides}
\end{align}

These quantities are evaluated continuously during Stage~2 powered flight: the
apogee radius $r_{\text{apo}}$ is compared with the target orbital radius in
the SECO event function (Section~\ref{sssec:seco_event}).

\subsubsection{Circularisation \texorpdfstring{$\Delta v$}{Delta-v}}
\label{sssec:circ_dv}

At apogee the velocity deficit with respect to a circular orbit is
\begin{equation}
  \Delta v_{\text{circ}} = \left|v_{\text{circ}} - v_{\text{apo}}\right|,
  \qquad v_{\text{circ}} = \sqrt{\frac{\mu}{r_{\text{target}}}},
  \label{eq:dv_circ}
\end{equation}
where $v_{\text{apo}}$ is the coast-phase velocity at apogee.  The propellant
mass required for this impulsive burn follows from the Tsiolkovsky equation:
\begin{equation}
  m_{\text{prop,circ}} = m_{\text{apo}}\!\left(1 - \exp\!\left[
    -\frac{\Delta v_{\text{circ}}}{I_{sp}\,g_0}\right]\right).
  \label{eq:mprop_circ}
\end{equation}

\subsubsection{Trajectory Losses and Gains}
\label{sssec:losses_gains}

The total velocity change delivered by the propulsion system exceeds the ideal
circular-orbit speed because several loss mechanisms are present during the
ascent:

\begin{itemize}
  \item \textbf{Gravity losses} arise because a component of gravity opposes the
        velocity during non-horizontal flight, so that not all thrust translates
        into kinetic-energy gain. These losses are most significant during the
        steep early ascent and are on the order of
        $1000$\,m\,s$^{-1}$ for a typical mission.

  \item \textbf{Drag losses} result from aerodynamic deceleration in the
        lower atmosphere, typically contributing ${\sim}100$\,m\,s$^{-1}$.

  \item \textbf{Steering losses} occur when the thrust vector is not aligned
        with the velocity direction ($\alpha \neq 0$), reducing the effective
        thrust along the trajectory.  Their magnitude depends on the guidance
        law and can reach several hundred m\,s$^{-1}$.
\end{itemize}

A positive contribution comes from \textbf{Earth rotation}: the launch site's
eastward surface velocity provides a free initial velocity increment
\begin{equation}
  \Delta v_{\text{rot}} = v_{\text{rot}}\sin\beta
  = \Omega_E\,R_E\cos\phi\;\sin\beta,
  \label{eq:dv_rot}
\end{equation}
where $\beta$ is the launch azimuth.  For an eastward launch from latitude
$28.5°$ this gain is approximately $410$\,m\,s$^{-1}$, a non-negligible
fraction of the orbital velocity budget.

% -----------------------------------------------------------------------------
\subsection{Guidance Laws}
\label{ssec:guidance}

Five guidance strategies are available in the simulator.  All share a common
initial pitchover manoeuvre; they differ in how the angle of attack~$\alpha$ is
computed after the vehicle exits the atmosphere.

\subsubsection{Common Initial Kick Phase}
\label{sssec:kick_phase}

Regardless of the guidance mode selected, the first seconds of flight follow
the same pitch programme.  The vehicle rises vertically until time $T_k$
(nominally $7.5$\,s), then executes a symmetric triangular pitchover manoeuvre
over a duration $\Delta T_k$ ($45$\,s):
\begin{equation}
  \alpha(t) =
  \begin{cases}
    0, & t < T_k, \\[2pt]
    \displaystyle\alpha_k\;\frac{2(t - T_k)}{\Delta T_k},
      & T_k \le t < T_k + \tfrac{\Delta T_k}{2}, \\[6pt]
    \displaystyle\alpha_k\;2\!\left(1 - \frac{t - T_k}{\Delta T_k}\right),
      & T_k + \tfrac{\Delta T_k}{2} \le t < T_k + \Delta T_k, \\[6pt]
    0, & t \ge T_k + \Delta T_k,
  \end{cases}
  \label{eq:kick_profile}
\end{equation}
where $\alpha_k$ is the \emph{initial kick angle}, the single design variable
optimised by the trajectory optimiser (Section~\ref{ssec:optimisation}).
Typical optimum values lie in the range $\alpha_k \in [-5°,\,-2°]$.

After the kick concludes and until the atmosphere-exit criterion is met,
$\alpha$ is held at zero.  Active guidance is then engaged.

\subsubsection{Gravity Turn}
\label{sssec:gt}

The gravity-turn mode applies no active steering after the initial kick:
\begin{equation}
  \alpha(t) = 0, \qquad \forall\; t > T_k + \Delta T_k.
  \label{eq:gt}
\end{equation}

The trajectory is shaped exclusively by gravity and the initial conditions.
This constitutes the simplest guidance baseline.

\subsubsection{Simple Polynomial}
\label{sssec:simple_poly}

Upon atmosphere exit, a linear steering law drives the flight-path angle toward
zero (horizontal flight for circular-orbit insertion):
\begin{equation}
  \alpha(\tau) = a_0 + a_1\,\tau,
  \label{eq:simple_poly}
\end{equation}
where $\tau = t_{\text{go}} / 100$ is a normalised time-to-go parameter, and
the coefficients are determined from the boundary conditions
\begin{equation}
  a_0 = \gamma_{\text{target}} = 0, \qquad
  a_1 = \gamma_{\text{current}}.
  \label{eq:simple_poly_bc}
\end{equation}

The time-to-go is estimated from the remaining altitude and the current radial
velocity component.  Coefficients are recomputed at a guidance update rate of
$0.5$\,s.

\subsubsection{Linear Tangent Steering}
\label{sssec:lts}

Linear tangent steering prescribes that the tangent of the total velocity
inclination angle $(\alpha + \gamma)$ varies linearly with time-to-go:
\begin{equation}
  \tan\bigl(\alpha + \gamma\bigr) = a\;t_{\text{go}} + b.
  \label{eq:lts}
\end{equation}

The boundary conditions are:
\begin{itemize}
  \item Terminal horizontal flight: $\tan(\alpha_f + \gamma_f) = 0$
        at $t_{\text{go}} = 0$, giving $b = 0$.
  \item Current-state continuity: at the present $t_{\text{go}}$,
        $\tan(\gamma) = a\;t_{\text{go}}$, giving
        $a = \tan\gamma / t_{\text{go}}$.
\end{itemize}

The commanded angle of attack is therefore
\begin{equation}
  \alpha = \arctan\!\bigl(a\;t_{\text{go}}\bigr) - \gamma.
  \label{eq:lts_alpha}
\end{equation}

This classical method provides smooth, monotonic inclination control and
is updated every $0.5$\,s.

\subsubsection{Bilinear Tangent Steering}
\label{sssec:bts}

The bilinear tangent law generalises linear tangent steering by representing
the inclination-angle tangent as a ratio of two linear functions of
$t_{\text{go}}$:
\begin{equation}
  \tan\bigl(\alpha + \gamma\bigr)
    = \frac{c_1\,t_{\text{go}} + c_2}
           {c_1'\,t_{\text{go}} + c_2'}.
  \label{eq:bts}
\end{equation}

The four coefficients are determined from four boundary conditions:
\begin{enumerate}
  \item \emph{Terminal value}: $c_2 = 0$,
        enforcing $\tan(\alpha_f + \gamma_f) = 0$ at $t_{\text{go}} = 0$.
  \item \emph{Terminal rate}: $c_1$ is set to a small constant
        (nominally $0.02$) to ensure a smooth, non-abrupt approach to
        horizontal flight.
  \item \emph{Normalisation}: $c_2' = 1$ removes the scaling ambiguity.
  \item \emph{Initial value}: $c_1'$ is computed so that the inclination
        at the current $t_{\text{go}}$ matches a blended target:
        \begin{equation}
          \gamma_{\text{blended}} = 0.7\;\gamma_{\text{current}}
            + 0.3\;\gamma_{\text{targeting}},
          \label{eq:gamma_blend}
        \end{equation}
        where $\gamma_{\text{targeting}}$ is a geometric estimate of the
        flight-path angle required to reach the target altitude over the
        remaining downrange distance.
\end{enumerate}

The bilinear form controls both the value and the rate of change of the
inclination angle at the terminal point, providing smoother orbit insertion
than the purely linear variant.

\subsubsection{Apollo Polynomial Guidance}
\label{sssec:apollo}

The most sophisticated guidance mode follows the explicit polynomial approach
originating from the Apollo programme.  Two independent channels command the
total acceleration in the downrange ($x$) and altitude ($y$) directions:
\begin{align}
  a_x(t) &= k_1\,(t - t_{\text{epoch}}) + k_2, \label{eq:apollo_ax} \\
  a_y(t) &= k_3\,(t - t_{\text{epoch}}) + k_4, \label{eq:apollo_ay}
\end{align}
where $t_{\text{epoch}}$ is the time at which the coefficients were last
recomputed.

\paragraph{Terminal constraints.}
The four coefficients $k_1,\dots,k_4$ are chosen to satisfy position and
velocity constraints at the target orbital altitude:
\begin{itemize}
  \item reach the target geocentric radius $r_f$,
  \item attain the local circular-orbit velocity
        $v_f = \sqrt{\mu / r_f}$ in the horizontal direction,
  \item achieve zero radial velocity ($\gamma_f = 0$).
\end{itemize}

The resulting expressions, derived from polynomial interpolation over the
interval $t_{\text{go}}$, are (shown for one generic channel with initial
state subscript~$i$ and final subscript~$f$):
\begin{align}
  k_1 &= \frac{6\,(v_f + v_i)\,t_{\text{go}}
          - 12\,(x_f - x_i)}{t_{\text{go}}^3},
        \label{eq:apollo_k1} \\[4pt]
  k_2 &= \frac{-2\,(v_f + 2\,v_i)\,t_{\text{go}}
          + 6\,(x_f - x_i)}{t_{\text{go}}^2}.
        \label{eq:apollo_k2}
\end{align}

Since the downrange target position is not known a~priori, it is estimated from
the current orbital elements by predicting where the unpowered trajectory would
intersect the target altitude.

\paragraph{Coefficient freezing.}
As $t_{\text{go}} \to 0$ the denominators $t_{\text{go}}^3$ and
$t_{\text{go}}^2$ drive the coefficients toward very large values.  To prevent
numerical instability the coefficients are frozen at their last computed values
once $t_{\text{go}}$ falls below a configurable threshold (nominally $10$\,s).

\paragraph{Steering-angle extraction.}
The total acceleration commanded by Eqs.~\eqref{eq:apollo_ax}--\eqref{eq:apollo_ay}
includes the gravitational contribution.  The thrust-only acceleration is
obtained by subtracting the local gravity vector, and the angle of attack is
then
\begin{equation}
  \alpha = \arctan\!\left(\frac{a_y^{\text{thrust}}}{a_x^{\text{thrust}}}\right)
           - \arctan\!\left(\frac{v_y}{v_x}\right).
  \label{eq:apollo_alpha}
\end{equation}

An optional thrust-magnitude-control mode allows the engine throttle to follow
the commanded acceleration magnitude
$|\mathbf{a}_{\text{thrust}}| = \sqrt{(a_x^{\text{thrust}})^2 +
(a_y^{\text{thrust}})^2}$, converting it to a thrust force via $F_T = m\,
|\mathbf{a}_{\text{thrust}}|$, limited by the maximum available thrust.

% -----------------------------------------------------------------------------
\subsection{Trajectory Optimisation}
\label{ssec:optimisation}

\subsubsection{Problem Formulation}
\label{sssec:opt_formulation}

The optimisation seeks the initial kick angle $\alpha_k$ that minimises the
total propellant consumption required to reach the target circular orbit:
\begin{equation}
  \min_{\alpha_k}\; J(\alpha_k)
    = m_{\text{prop,ascent}}(\alpha_k)
    + m_{\text{prop,circ}}(\alpha_k).
  \label{eq:obj}
\end{equation}

The single design variable is bounded:
$\alpha_k \in [\alpha_{\min},\,\alpha_{\max}]$ (nominally $[-5°,\,-2°]$).

Three soft constraints are enforced through a large penalty value
($J = 10^9$\,kg) applied when any condition is violated:
\begin{enumerate}
  \item propellant required for circularisation exceeds the available supply,
  \item the circularisation burn duration exceeds a maximum limit
        (nominally $15$\,s),
  \item the achieved apogee deviates from the target by more than $0.2\,\%$.
\end{enumerate}

\subsubsection{Optimisation Algorithm}
\label{sssec:opt_algorithm}

Because the objective function is non-smooth --- discrete events such as
staging, engine ignitions, and cutoffs introduce discontinuities --- a
gradient-free approach is used.  A brute-force grid search evaluates
$J(\alpha_k)$ at $N = 1000$ uniformly spaced points across
$[\alpha_{\min},\,\alpha_{\max}]$ and selects the global minimum.  No local
refinement step is applied.

The grid resolution is approximately $0.003°$, which is sufficient for
engineering accuracy.  With each trajectory evaluation taking on the order of
one second, the full optimisation completes in roughly $15$--$20$~minutes.

A key architectural feature is that the optimisation loop is entirely
independent of the guidance mode.  The outer loop supplies a candidate kick
angle and receives back a scalar cost; the inner simulation selects and
executes whichever guidance law is active.  Consequently the same search
procedure applies unchanged to all five guidance modes.

\subsubsection{Event-Based Engine Cutoff}
\label{sssec:seco_event}

The timing of SECO is not a free optimisation variable.  Instead, it is
determined implicitly during trajectory integration through an event function:
\begin{equation}
  g(t,\mathbf{x}) = r_{\text{apo}}(t,\mathbf{x}) - r_{\text{target}},
  \label{eq:seco_event}
\end{equation}
where $r_{\text{apo}}$ is the instantaneous apogee computed from the current
orbital elements (Eq.~\ref{eq:apsides}).  The integrator monitors $g$ during
the Stage~2 burn (above $65$\,km altitude); a sign change triggers root-finding
to locate the exact cutoff time.

The physical interpretation is direct: the engine burns until the trajectory,
if left unpowered from that instant onward, would coast to exactly the target
apogee altitude.  This event-driven approach automatically adapts the burn
duration to the trajectory shape produced by each guidance law and kick angle.

% -----------------------------------------------------------------------------
\subsection{Space Launcher Configuration}
\label{ssec:launcher}

The launch vehicle modelled in the simulator is representative of a
two-stage Falcon~9-class rocket.  Table~\ref{tab:launcher} summarises
the key parameters.

\begin{table}[htbp]
  \centering
  \caption{Launch-vehicle parameters used in the simulator.}
  \label{tab:launcher}
  \begin{tabular}{l c c c}
    \toprule
    \textbf{Parameter} & \textbf{Stage~1} & \textbf{Stage~2} & \textbf{Unit} \\
    \midrule
    Thrust (vacuum)       & 7\,600  & 934    & kN   \\
    Specific impulse      & 283     & 348    & s    \\
    Structural mass       & 25\,600 & 3\,900 & kg   \\
    Propellant mass       & 395\,700& 92\,670& kg   \\
    \bottomrule
  \end{tabular}
\end{table}

Stage separation occurs $3$\,s after MECO, and the second-stage engine ignites
$8$\,s after MECO, giving a $5$\,s coast interval between stages.

The aerodynamic reference area is $A = 10.52$\,m$^2$ with a constant drag
coefficient $C_D = 0.3$.

\subsubsection{Mission Parameters}

The baseline mission targets a $500$\,km circular orbit at an inclination of
$51.6°$ (representative of the International Space Station orbit), launching
from a latitude of $28.5°$ (Kennedy Space Center).  The resulting initial mass
at lift-off is approximately $517.9$\,t (sum of both stages, with zero payload
in the reference configuration).  Table~\ref{tab:mission} lists the nominal
mission parameters.

\begin{table}[htbp]
  \centering
  \caption{Baseline mission parameters.}
  \label{tab:mission}
  \begin{tabular}{l c c}
    \toprule
    \textbf{Parameter} & \textbf{Value} & \textbf{Unit} \\
    \midrule
    Target orbital altitude   & 500    & km   \\
    Target inclination        & 51.6   & deg  \\
    Launch latitude           & 28.5   & deg  \\
    Kick-angle search range   & $[-5,\,-2]$ & deg \\
    Guidance update rate      & 0.5    & s    \\
    Integration tolerance     & $10^{-8}$ & --- \\
    Maximum integration step  & 1.0    & s    \\
    Output time resolution    & 0.01   & s    \\
    \bottomrule
  \end{tabular}
\end{table}

% =============================================================================

\section{Simulator Implementation}
\label{sec:simulator_implementation}

This section describes how the trajectory simulation tool is structured and how
the theoretical principles introduced in the previous chapter are brought
together into a working ascent simulator.  The discussion covers the overall
software architecture, the reference frames adopted, the equations of motion
actually integrated, the environmental and force models, the guidance laws
available, and the optimisation strategy used to find the best initial
conditions.  Vehicle parameters representative of a Falcon~9-class launcher
close the section.

% -----------------------------------------------------------------------------
\subsection{Simulation Architecture and Workflow}
\label{ssec:sim_architecture}

The simulator is organised around three successive mission phases, each
resolved by a separate call to the numerical integrator:

\begin{enumerate}
  \item \textbf{Stage~1 powered ascent.}  From lift-off until first-stage
        propellant depletion (Main Engine Cut-Off, MECO).  During this phase
        the vehicle follows an initial vertical rise, a pitchover kick
        manoeuvre, and atmospheric flight with aerodynamic drag.

  \item \textbf{Stage~2 powered ascent with guidance.}  After stage separation
        and a brief coast interval, the second-stage engine ignites.  An active
        guidance law (selected from five available modes) shapes the trajectory
        until the event-triggered Second Engine Cut-Off (SECO), which occurs
        when the instantaneous apogee of the coasting trajectory matches the
        target orbital altitude.

  \item \textbf{Ballistic coast and circularisation.}  The vehicle coasts
        unpowered to apogee, where an impulsive circularisation burn is applied
        to achieve the target circular orbit.  A final short coast in stable
        orbit concludes the simulation.
\end{enumerate}

\subsubsection{State Vector}

The base state vector integrated by the solver is
\begin{equation}
  \mathbf{x} = \bigl[\,s,\; r,\; v,\; \gamma,\; m\,\bigr]^T,
  \label{eq:state_base}
\end{equation}
where $s$ is the downrange distance along the Earth surface, $r$ is the
geocentric radius, $v$ is the velocity magnitude in the rotating frame,
$\gamma$ is the flight-path angle, and $m$ is the vehicle mass.

When Earth-rotation effects are enabled, the state is extended to include
geocentric latitude~$\phi$:
\begin{equation}
  \mathbf{x} = \bigl[\,s,\; r,\; v,\; \gamma,\; m,\; \phi\,\bigr]^T.
  \label{eq:state_rot}
\end{equation}

If heading tracking is additionally activated, the local heading (azimuth)
$\psi$ becomes the seventh state component:
\begin{equation}
  \mathbf{x} = \bigl[\,s,\; r,\; v,\; \gamma,\; m,\; \phi,\; \psi\,\bigr]^T.
  \label{eq:state_heading}
\end{equation}

\subsubsection{Numerical Integration}

Each mission phase is propagated with an explicit Runge--Kutta~4(5) scheme
(Dormand--Prince variant) provided by \texttt{scipy.integrate.solve\_ivp}.  The
solver uses adaptive step-size control with an absolute tolerance of
$10^{-8}$, a maximum integration step of $1$~s, and a dense-output time
resolution of $0.01$~s.

Phase transitions are handled through \emph{event functions}: scalar functions
of the state that the integrator monitors for zero-crossings.  When a crossing
is detected the integrator performs root-finding to locate the exact transition
time and terminates the current phase.  Events used include propellant
depletion (MECO), stage separation timing, atmosphere exit, target-apogee
matching (SECO), and ground-collision safeguards.

\subsubsection{Architecture Overview}

Figure~\ref{fig:sim_architecture} shows the high-level data flow.  An outer
optimisation loop varies the initial kick angle supplied to the simulation
engine.  Inside each simulation call, the equations of motion are evaluated at
every integrator step using the current environmental models (gravity,
atmosphere) and the active guidance law to compute the steering command.  Event
functions run concurrently and trigger phase switches when their conditions are
met.  After the full trajectory is obtained, the optimiser evaluates the
objective function and proceeds to the next candidate.

\begin{figure}[htbp]
  \centering
  \begin{tikzpicture}[
      >=latex,
      block/.style  = {draw, rounded corners, minimum height=1.0cm,
                        minimum width=2.6cm, align=center, font=\small},
      smallblock/.style = {draw, rounded corners, minimum height=0.8cm,
                        minimum width=2.2cm, align=center, font=\footnotesize},
      arrow/.style  = {->, thick},
      darrow/.style = {<->, thick},
    ]

    % Optimisation loop
    \node[block, fill=blue!12] (opt) {Optimisation\\Loop};

    % Simulation engine
    \node[block, fill=green!10, right=2.2cm of opt] (sim) {Simulation\\Engine};

    % Inner modules (below sim)
    \node[smallblock, fill=orange!10, below left=1.4cm and -0.3cm of sim]
          (guid) {Guidance\\Module};
    \node[smallblock, fill=orange!10, below=1.4cm of sim]
          (eom) {Equations\\of Motion};
    \node[smallblock, fill=orange!10, below right=1.4cm and -0.3cm of sim]
          (env) {Environment\\Models};

    % Event detection
    \node[block, fill=red!10, right=2.2cm of sim] (evt) {Event\\Detection};

    % Outputs
    \node[block, fill=gray!12, right=2.2cm of evt] (out) {Trajectory\\Outputs};

    % Arrows
    \draw[arrow] (opt) -- node[above, font=\footnotesize]{$\alpha_{\text{kick}}$} (sim);
    \draw[arrow] (sim) -- (evt);
    \draw[arrow] (evt) -- (out);
    \draw[arrow] (out.south) -- ++(0,-1.2) -| node[near start, below, font=\footnotesize]{$J(\alpha)$} (opt.south);

    % Inner connections
    \draw[arrow] (sim) -- (guid);
    \draw[arrow] (sim) -- (eom);
    \draw[arrow] (sim) -- (env);
    \draw[arrow] (guid) -- (eom);
    \draw[arrow] (env) -- (eom);

  \end{tikzpicture}
  \caption{High-level architecture of the trajectory simulator.  The
           optimisation loop feeds candidate kick angles into the simulation
           engine, which dispatches environment evaluations and guidance
           commands to the equations of motion at each integration step.  Event
           detection triggers phase transitions and engine cutoff.  The
           resulting trajectory is scored and fed back to the optimiser.}
  \label{fig:sim_architecture}
\end{figure}

% -----------------------------------------------------------------------------
\subsection{Reference Frames}
\label{ssec:reference_frames}

The simulator employs a mixed-frame strategy that keeps guidance and
atmosphere coupling straightforward while preserving correct orbital-mechanics
metrics.

\begin{itemize}
  \item \textbf{Rotating frame (ECEF-like).}  All ascent dynamics quantities
        --- velocity magnitude~$v$, flight-path angle~$\gamma$, and heading
        decomposition --- are expressed in a frame that co-rotates with the
        Earth.  This simplifies the treatment of aerodynamic drag, which
        depends on the velocity relative to the atmosphere.

  \item \textbf{Inertial frame (ECI).}  Orbital-energy and orbital-element
        calculations (semi-major axis, eccentricity, apogee radius) are
        performed in an Earth-Centred Inertial frame.  This is essential for
        correct event detection (e.g.\ target-apogee matching) and end-of-run
        orbit characterisation.

  \item \textbf{Local East--North--Up (ENU).}  Pseudo-force accelerations
        (Coriolis and centrifugal) are first evaluated in the local ENU basis at
        the vehicle's current latitude, then projected onto the along-heading
        and radial directions used by the equations of motion.
\end{itemize}

\subsubsection{Frame Transition During Flight}

During powered ascent and early coasting the integration proceeds in the
rotating frame with optional pseudo-force corrections.  Once the vehicle
reaches the final ballistic coast phase (post-SECO), the propagation switches
to pure inertial dynamics: pseudo-force terms are deactivated to avoid
introducing spurious accelerations when the trajectory is governed solely by
two-body gravity.

\subsubsection{ECEF-to-ECI Velocity Conversion}

At any instant the local rotating-frame velocity components are converted to
their inertial counterparts by adding the Earth's rotational eastward speed.
The horizontal velocity magnitude is $v_h = v\cos\gamma$, which decomposes as
\begin{equation}
  v_N = v_h\cos\psi, \qquad
  v_E^{\text{ECEF}} = v_h\sin\psi.
  \label{eq:vel_ecef}
\end{equation}

Adding Earth rotation yields
\begin{equation}
  v_E^{\text{ECI}} = v_E^{\text{ECEF}} + \Omega_E\, r\cos\phi,
  \label{eq:ve_eci}
\end{equation}
from which the inertial speed and flight-path angle follow:
\begin{align}
  v_{\text{ECI}} &= \sqrt{\bigl(v_E^{\text{ECI}}\bigr)^2 + v_N^2
                    + \bigl(v\sin\gamma\bigr)^2}, \label{eq:v_eci} \\[4pt]
  \gamma_{\text{ECI}} &= \operatorname{atan2}\!\bigl(v\sin\gamma,\;
                    \sqrt{\bigl(v_E^{\text{ECI}}\bigr)^2 + v_N^2}\;\bigr).
                    \label{eq:gamma_eci}
\end{align}

These inertial quantities are used whenever orbital elements are computed.

% -----------------------------------------------------------------------------
\subsection{Launch Azimuth and Heading}
\label{ssec:azimuth_heading}

\subsubsection{Launch Azimuth Determination}
\label{sssec:launch_azimuth}

The launch azimuth defines the initial heading of the ascent trajectory in the
horizontal plane and directly controls the achievable orbital inclination.
Two modes are available.

\paragraph{Geometric (inertial) azimuth.}
Given a target inclination~$i$ and launch latitude~$\phi$, the classical
spherical-geometry relation gives the inertial azimuth measured from true
north:
\begin{equation}
  \sin\beta_{\text{geom}} = \frac{\cos i}{\cos\phi}.
  \label{eq:az_geom}
\end{equation}

\paragraph{Corrected (rotating-frame) azimuth.}
Because the rotating-frame velocity seen by the guidance system differs from
the inertial velocity by the eastward surface speed, the launch azimuth is
corrected:
\begin{equation}
  v_{\text{rot}} = \Omega_E\, R_E\cos\phi,
  \label{eq:vrot}
\end{equation}
\begin{align}
  v_N &= v_{\text{orb}}\cos\beta_{\text{geom}}, \label{eq:vn_corr} \\
  v_E^{\text{ECEF}} &= v_{\text{orb}}\sin\beta_{\text{geom}} - v_{\text{rot}},
  \label{eq:ve_corr}
\end{align}
where $v_{\text{orb}} = \sqrt{\mu/r_{\text{target}}}$ is the circular speed at
the target altitude.  The corrected azimuth is then
\begin{equation}
  \beta_{\text{corr}} = \operatorname{atan2}\!\bigl(v_E^{\text{ECEF}},\, v_N\bigr).
  \label{eq:az_corr}
\end{equation}

The simulator defaults to the corrected mode; a configuration flag may select
the geometric mode instead.

\subsubsection{Heading Propagation}
\label{sssec:heading_propagation}

When heading tracking is enabled, the local azimuth $\psi$ is propagated as a
state variable.  Its rate is computed from spherical-geometry consistency with
the great-circle ground track:
\begin{equation}
  \dot\psi = \tan\phi\;\tan\psi\;\dot\phi,
  \label{eq:psi_dot}
\end{equation}
where $\dot\phi$ is obtained from the latitude kinematic relation detailed in
Section~\ref{ssec:eom}.  An optional cross-heading pseudo-force contribution
may be added:
\begin{equation}
  \dot\psi \mathrel{+}= \frac{a_{\perp}}{v\cos\gamma},
  \label{eq:psi_pseudo}
\end{equation}
with $a_{\perp}$ the component of the Coriolis and centrifugal acceleration
perpendicular to the heading direction (see Section~\ref{sssec:pseudo_forces}).

\subsubsection{Achieved Inclination}
\label{sssec:achieved_inclination}

At the end of the simulation the achieved orbital inclination is evaluated by
converting the final rotating-frame velocity to inertial components using
Eqs.~\eqref{eq:vel_ecef}--\eqref{eq:ve_eci} and computing the inertial
velocity azimuth:
\begin{equation}
  \beta_{\text{ECI}} = \operatorname{atan2}\!\bigl(v_E^{\text{ECI}},\, v_N\bigr).
  \label{eq:beta_eci}
\end{equation}

The inclination follows from the classical relation:
\begin{equation}
  \cos i_{\text{achieved}} = \cos\phi\;\sin\beta_{\text{ECI}}.
  \label{eq:inc_achieved}
\end{equation}

The inclination drift $\Delta i = i_{\text{achieved}} - i_{\text{target}}$ is
a valuable metric for assessing the fidelity of the rotating-frame treatment
and the effectiveness of the heading-tracking augmentation.

% -----------------------------------------------------------------------------
\subsection{Equations of Motion}
\label{ssec:eom}

\subsubsection{Base Dynamics (Non-Rotating Form)}

The core differential system, expressed in spherical coordinates before
optional pseudo-force augmentation, is
\begin{align}
  \dot s      &= \frac{R_E}{r}\,v\cos\gamma,
               \label{eq:sdot}\\[3pt]
  \dot r      &= v\sin\gamma,
               \label{eq:rdot}\\[3pt]
  \dot v      &= \frac{F_T}{m}\cos\alpha
                 - \frac{F_D}{m}
                 - g(r)\sin\gamma,
               \label{eq:vdot}\\[3pt]
  \dot\gamma  &= \frac{1}{v}\!\left(
                   \frac{F_T}{m}\sin\alpha
                   + \frac{F_L}{m}
                   - \Bigl[g(r) - \frac{v^2}{r}\Bigr]\cos\gamma
                 \right),
               \label{eq:gammadot}\\[3pt]
  \dot m      &= -\frac{F_T}{I_{sp}\,g_0},
               \label{eq:mdot}
\end{align}
where $\alpha$ is the angle of attack (steering command from the guidance law),
$F_T$ and $I_{sp}$ are the current thrust and specific impulse, $F_D$ is
aerodynamic drag, $F_L$ is lift (set to zero in the present implementation),
and $g(r) = \mu/r^2$ is the gravitational acceleration.

The term $v^2/r$ in Eq.~\eqref{eq:gammadot} represents the centripetal
acceleration arising from the curved trajectory; when $v^2/r$ exceeds
$g\cos\gamma$ the net effect tends to raise the flight-path angle.

Two numerical safeguards are applied:
\begin{itemize}
  \item \emph{Small-velocity guard:} when $v < 10^{-6}$\,m/s the angular rate
        $\dot\gamma$ is forced to zero to prevent a singularity.
  \item \emph{Vertical-hold constraint:} during the initial vertical rise
        (before the pitchover kick begins) $\dot\gamma$ is explicitly set to
        zero regardless of geometric terms.
\end{itemize}

\subsubsection{Rotating-Frame Augmentation}

When Earth-rotation effects are enabled, Coriolis and centrifugal
accelerations computed in the local ENU frame (see
Section~\ref{sssec:pseudo_forces}) are projected onto the along-heading and
radial directions:
\begin{equation}
  a_{h,\parallel} = a_E\sin\psi + a_N\cos\psi,
  \label{eq:a_parallel}
\end{equation}
where $a_E$ and $a_N$ collect both Coriolis and centrifugal components.
The additive corrections to the velocity and flight-path-angle rates are
\begin{align}
  \Delta\dot v     &= a_{h,\parallel}\cos\gamma + a_{\text{up}}\sin\gamma,
                    \label{eq:dv_pseudo}\\[3pt]
  \Delta\dot\gamma &= \frac{-a_{h,\parallel}\sin\gamma
                      + a_{\text{up}}\cos\gamma}{v},
                    \label{eq:dgamma_pseudo}
\end{align}
so that the integrated rates become
$\dot v_{\text{tot}} = \dot v + \Delta\dot v$ and
$\dot\gamma_{\text{tot}} = \dot\gamma + \Delta\dot\gamma$.

\subsubsection{Latitude Propagation}

Rather than integrating a differential latitude rate directly, the simulator
reconstructs latitude from the downrange angle using great-circle geometry,
which avoids numerical drift:
\begin{equation}
  \sigma = \frac{s}{R_E}, \qquad
  \sin\phi = \sin\phi_0\cos\sigma + \cos\phi_0\sin\sigma\cos\beta_0,
  \label{eq:lat_recon}
\end{equation}
where $\phi_0$ is the launch latitude and $\beta_0$ is the inertial launch
azimuth.  The latitude rate required by the heading equation
\eqref{eq:psi_dot} is obtained by differentiating analytically:
\begin{equation}
  \dot\phi = \frac{-\sin\phi_0\sin\sigma
    + \cos\phi_0\cos\sigma\cos\beta_0}{\cos\phi}\;\frac{\dot s}{R_E}.
  \label{eq:phidot}
\end{equation}

% -----------------------------------------------------------------------------
\subsection{Environmental Models}
\label{ssec:env_models}

\subsubsection{Gravitational Model}
\label{sssec:gravity}

Gravity is modelled as a central inverse-square field:
\begin{equation}
  g(r) = \frac{\mu}{r^2},
  \label{eq:gravity}
\end{equation}
with $\mu = 3.986\times10^{14}$\,m$^3$\,s$^{-2}$.  The Earth is assumed
spherical; higher-order harmonics such as $J_2$ are not included.

\subsubsection{Atmospheric Model}
\label{sssec:atmosphere}

Atmospheric density follows an exponential profile:
\begin{equation}
  \rho(h) = \rho_0\,\exp\!\left(-\frac{h}{H}\right),
  \qquad h = r - R_E,
  \label{eq:atmosphere}
\end{equation}
with sea-level density $\rho_0 = 1.225$\,kg\,m$^{-3}$ and scale height
$H = 8500$\,m.  The dynamic pressure is $q = \tfrac{1}{2}\rho\,v^2$.

Two criteria are available for declaring atmosphere exit and
activating exo-atmospheric guidance:
\begin{itemize}
  \item \emph{Altitude threshold:} $h > 65$\,km.
  \item \emph{Dynamic-pressure threshold:} $q < 1000$\,Pa (configurable).
\end{itemize}
The selected criterion is a simulation parameter.

% -----------------------------------------------------------------------------
\subsection{External Forces}
\label{ssec:forces}

\subsubsection{Aerodynamic Forces}
\label{sssec:aero}

The aerodynamic drag force opposing the velocity vector is
\begin{equation}
  F_D = q\,C_D\,A = \tfrac{1}{2}\,\rho\,v^2\,C_D\,A,
  \label{eq:drag}
\end{equation}
with constant drag coefficient $C_D = 0.3$ and reference area $A = 10.52$\,m$^2$.
A lift model is available in the formulation (Eq.~\ref{eq:gammadot}), but in
all analyses presented here lift is neglected ($F_L = 0$).

\subsubsection{Propulsive Forces}
\label{sssec:propulsion}

Thrust magnitude $F_T$ and specific impulse $I_{sp}$ are stage-dependent
constants that change at discrete events (MECO, stage separation, second-stage
ignition, SECO).  Thrust is directed at an angle of attack $\alpha$ relative
to the velocity vector; $\alpha$ is the primary control variable and is
determined at each integration step by the active guidance law
(Section~\ref{ssec:guidance}).

The resulting vehicle mass depletion obeys
\begin{equation}
  \dot m = -\frac{F_T}{I_{sp}\,g_0},
  \label{eq:mdot_prop}
\end{equation}
where $g_0 = 9.80665$\,m\,s$^{-2}$ is the standard gravitational acceleration.

\subsubsection{Pseudo-Forces (Earth Rotation)}
\label{sssec:pseudo_forces}

When the integration is performed in the rotating frame with Earth-rotation
corrections enabled, inertial pseudo-accelerations are applied.

\paragraph{Coriolis acceleration.}
In the local ENU basis at latitude $\phi$:
\begin{align}
  a_E^{\text{cor}} &= -2\Omega_E\cos\phi\;v_{\text{up}}
                      + 2\Omega_E\sin\phi\;v_N,
                    \label{eq:cor_e} \\
  a_N^{\text{cor}} &= -2\Omega_E\sin\phi\;v_E,
                    \label{eq:cor_n} \\
  a_{\text{up}}^{\text{cor}} &= 2\Omega_E\cos\phi\;v_E.
                    \label{eq:cor_u}
\end{align}

\paragraph{Centrifugal acceleration.}
\begin{align}
  a_E^{\text{cen}} &= 0,
                    \label{eq:cen_e} \\
  a_N^{\text{cen}} &= -\Omega_E^2\,r\,\sin\phi\cos\phi,
                    \label{eq:cen_n} \\
  a_{\text{up}}^{\text{cen}} &= \Omega_E^2\,r\,\cos^2\!\phi.
                    \label{eq:cen_u}
\end{align}

\paragraph{Projection onto flight-dynamics directions.}
The total pseudo-acceleration in each ENU direction is summed ($a_E = a_E^{\text{cor}} + a_E^{\text{cen}}$, etc.) and then projected onto the heading direction and the vertical:
\begin{align}
  a_{h,\parallel} &= a_E\sin\psi + a_N\cos\psi
    & &\text{(along-heading)}, \label{eq:a_par} \\
  a_{h,\perp}     &= a_E\cos\psi - a_N\sin\psi
    & &\text{(cross-heading)}, \label{eq:a_perp}
\end{align}
with $a_{\text{up}}$ retained for the radial direction.

The along-heading and up components enter the velocity and flight-path-angle
corrections (Eqs.~\ref{eq:dv_pseudo}--\ref{eq:dgamma_pseudo}).  If cross-heading
pseudo-force tracking is enabled, $a_{h,\perp}$ contributes to the heading
rate via Eq.~\eqref{eq:psi_pseudo}.

% -----------------------------------------------------------------------------
\subsection{Performance Metrics}
\label{ssec:performance}

\subsubsection{Orbital Elements}
\label{sssec:orbital_elements}

Once the velocity is converted to the inertial frame
(Eqs.~\ref{eq:v_eci}--\ref{eq:gamma_eci}), the standard two-body orbital
elements are computed:
\begin{align}
  a &= \frac{\mu\,r}{2\mu - r\,v_{\text{ECI}}^2},
    \label{eq:sma}\\[3pt]
  e &= \sqrt{1 - \frac{\bigl(r\,v_{\text{ECI}}\cos\gamma_{\text{ECI}}\bigr)^2}
       {\mu\,a}},
    \label{eq:ecc}\\[3pt]
  r_{\text{apo}} &= a\,(1+e), \qquad
  r_{\text{peri}} = a\,(1-e).
  \label{eq:apsides}
\end{align}

These quantities are evaluated continuously during Stage~2 powered flight: the
apogee radius $r_{\text{apo}}$ is compared with the target orbital radius in
the SECO event function (Section~\ref{sssec:seco_event}).

\subsubsection{Circularisation \texorpdfstring{$\Delta v$}{Delta-v}}
\label{sssec:circ_dv}

At apogee the velocity deficit with respect to a circular orbit is
\begin{equation}
  \Delta v_{\text{circ}} = \left|v_{\text{circ}} - v_{\text{apo}}\right|,
  \qquad v_{\text{circ}} = \sqrt{\frac{\mu}{r_{\text{target}}}},
  \label{eq:dv_circ}
\end{equation}
where $v_{\text{apo}}$ is the coast-phase velocity at apogee.  The propellant
mass required for this impulsive burn follows from the Tsiolkovsky equation:
\begin{equation}
  m_{\text{prop,circ}} = m_{\text{apo}}\!\left(1 - \exp\!\left[
    -\frac{\Delta v_{\text{circ}}}{I_{sp}\,g_0}\right]\right).
  \label{eq:mprop_circ}
\end{equation}

\subsubsection{Trajectory Losses and Gains}
\label{sssec:losses_gains}

The total velocity change delivered by the propulsion system exceeds the ideal
circular-orbit speed because several loss mechanisms are present during the
ascent:

\begin{itemize}
  \item \textbf{Gravity losses} arise because a component of gravity opposes the
        velocity during non-horizontal flight, so that not all thrust translates
        into kinetic-energy gain. These losses are most significant during the
        steep early ascent and are on the order of
        $1000$\,m\,s$^{-1}$ for a typical mission.

  \item \textbf{Drag losses} result from aerodynamic deceleration in the
        lower atmosphere, typically contributing ${\sim}100$\,m\,s$^{-1}$.

  \item \textbf{Steering losses} occur when the thrust vector is not aligned
        with the velocity direction ($\alpha \neq 0$), reducing the effective
        thrust along the trajectory.  Their magnitude depends on the guidance
        law and can reach several hundred m\,s$^{-1}$.
\end{itemize}

A positive contribution comes from \textbf{Earth rotation}: the launch site's
eastward surface velocity provides a free initial velocity increment
\begin{equation}
  \Delta v_{\text{rot}} = v_{\text{rot}}\sin\beta
  = \Omega_E\,R_E\cos\phi\;\sin\beta,
  \label{eq:dv_rot}
\end{equation}
where $\beta$ is the launch azimuth.  For an eastward launch from latitude
$28.5°$ this gain is approximately $410$\,m\,s$^{-1}$, a non-negligible
fraction of the orbital velocity budget.

% -----------------------------------------------------------------------------
\subsection{Guidance Laws}
\label{ssec:guidance}

Five guidance strategies are available in the simulator.  All share a common
initial pitchover manoeuvre; they differ in how the angle of attack~$\alpha$ is
computed after the vehicle exits the atmosphere.

\subsubsection{Common Initial Kick Phase}
\label{sssec:kick_phase}

Regardless of the guidance mode selected, the first seconds of flight follow
the same pitch programme.  The vehicle rises vertically until time $T_k$
(nominally $7.5$\,s), then executes a symmetric triangular pitchover manoeuvre
over a duration $\Delta T_k$ ($45$\,s):
\begin{equation}
  \alpha(t) =
  \begin{cases}
    0, & t < T_k, \\[2pt]
    \displaystyle\alpha_k\;\frac{2(t - T_k)}{\Delta T_k},
      & T_k \le t < T_k + \tfrac{\Delta T_k}{2}, \\[6pt]
    \displaystyle\alpha_k\;2\!\left(1 - \frac{t - T_k}{\Delta T_k}\right),
      & T_k + \tfrac{\Delta T_k}{2} \le t < T_k + \Delta T_k, \\[6pt]
    0, & t \ge T_k + \Delta T_k,
  \end{cases}
  \label{eq:kick_profile}
\end{equation}
where $\alpha_k$ is the \emph{initial kick angle}, the single design variable
optimised by the trajectory optimiser (Section~\ref{ssec:optimisation}).
Typical optimum values lie in the range $\alpha_k \in [-5°,\,-2°]$.

After the kick concludes and until the atmosphere-exit criterion is met,
$\alpha$ is held at zero.  Active guidance is then engaged.

\subsubsection{Gravity Turn}
\label{sssec:gt}

The gravity-turn mode applies no active steering after the initial kick:
\begin{equation}
  \alpha(t) = 0, \qquad \forall\; t > T_k + \Delta T_k.
  \label{eq:gt}
\end{equation}

The trajectory is shaped exclusively by gravity and the initial conditions.
This constitutes the simplest guidance baseline.

\subsubsection{Simple Polynomial}
\label{sssec:simple_poly}

Upon atmosphere exit, a linear steering law drives the flight-path angle toward
zero (horizontal flight for circular-orbit insertion):
\begin{equation}
  \alpha(\tau) = a_0 + a_1\,\tau,
  \label{eq:simple_poly}
\end{equation}
where $\tau = t_{\text{go}} / 100$ is a normalised time-to-go parameter, and
the coefficients are determined from the boundary conditions
\begin{equation}
  a_0 = \gamma_{\text{target}} = 0, \qquad
  a_1 = \gamma_{\text{current}}.
  \label{eq:simple_poly_bc}
\end{equation}

The time-to-go is estimated from the remaining altitude and the current radial
velocity component.  Coefficients are recomputed at a guidance update rate of
$0.5$\,s.

\subsubsection{Linear Tangent Steering}
\label{sssec:lts}

Linear tangent steering prescribes that the tangent of the total velocity
inclination angle $(\alpha + \gamma)$ varies linearly with time-to-go:
\begin{equation}
  \tan\bigl(\alpha + \gamma\bigr) = a\;t_{\text{go}} + b.
  \label{eq:lts}
\end{equation}

The boundary conditions are:
\begin{itemize}
  \item Terminal horizontal flight: $\tan(\alpha_f + \gamma_f) = 0$
        at $t_{\text{go}} = 0$, giving $b = 0$.
  \item Current-state continuity: at the present $t_{\text{go}}$,
        $\tan(\gamma) = a\;t_{\text{go}}$, giving
        $a = \tan\gamma / t_{\text{go}}$.
\end{itemize}

The commanded angle of attack is therefore
\begin{equation}
  \alpha = \arctan\!\bigl(a\;t_{\text{go}}\bigr) - \gamma.
  \label{eq:lts_alpha}
\end{equation}

This classical method provides smooth, monotonic inclination control and
is updated every $0.5$\,s.

\subsubsection{Bilinear Tangent Steering}
\label{sssec:bts}

The bilinear tangent law generalises linear tangent steering by representing
the inclination-angle tangent as a ratio of two linear functions of
$t_{\text{go}}$:
\begin{equation}
  \tan\bigl(\alpha + \gamma\bigr)
    = \frac{c_1\,t_{\text{go}} + c_2}
           {c_1'\,t_{\text{go}} + c_2'}.
  \label{eq:bts}
\end{equation}

The four coefficients are determined from four boundary conditions:
\begin{enumerate}
  \item \emph{Terminal value}: $c_2 = 0$,
        enforcing $\tan(\alpha_f + \gamma_f) = 0$ at $t_{\text{go}} = 0$.
  \item \emph{Terminal rate}: $c_1$ is set to a small constant
        (nominally $0.02$) to ensure a smooth, non-abrupt approach to
        horizontal flight.
  \item \emph{Normalisation}: $c_2' = 1$ removes the scaling ambiguity.
  \item \emph{Initial value}: $c_1'$ is computed so that the inclination
        at the current $t_{\text{go}}$ matches a blended target:
        \begin{equation}
          \gamma_{\text{blended}} = 0.7\;\gamma_{\text{current}}
            + 0.3\;\gamma_{\text{targeting}},
          \label{eq:gamma_blend}
        \end{equation}
        where $\gamma_{\text{targeting}}$ is a geometric estimate of the
        flight-path angle required to reach the target altitude over the
        remaining downrange distance.
\end{enumerate}

The bilinear form controls both the value and the rate of change of the
inclination angle at the terminal point, providing smoother orbit insertion
than the purely linear variant.

\subsubsection{Apollo Polynomial Guidance}
\label{sssec:apollo}

The most sophisticated guidance mode follows the explicit polynomial approach
originating from the Apollo programme.  Two independent channels command the
total acceleration in the downrange ($x$) and altitude ($y$) directions:
\begin{align}
  a_x(t) &= k_1\,(t - t_{\text{epoch}}) + k_2, \label{eq:apollo_ax} \\
  a_y(t) &= k_3\,(t - t_{\text{epoch}}) + k_4, \label{eq:apollo_ay}
\end{align}
where $t_{\text{epoch}}$ is the time at which the coefficients were last
recomputed.

\paragraph{Terminal constraints.}
The four coefficients $k_1,\dots,k_4$ are chosen to satisfy position and
velocity constraints at the target orbital altitude:
\begin{itemize}
  \item reach the target geocentric radius $r_f$,
  \item attain the local circular-orbit velocity
        $v_f = \sqrt{\mu / r_f}$ in the horizontal direction,
  \item achieve zero radial velocity ($\gamma_f = 0$).
\end{itemize}

The resulting expressions, derived from polynomial interpolation over the
interval $t_{\text{go}}$, are (shown for one generic channel with initial
state subscript~$i$ and final subscript~$f$):
\begin{align}
  k_1 &= \frac{6\,(v_f + v_i)\,t_{\text{go}}
          - 12\,(x_f - x_i)}{t_{\text{go}}^3},
        \label{eq:apollo_k1} \\[4pt]
  k_2 &= \frac{-2\,(v_f + 2\,v_i)\,t_{\text{go}}
          + 6\,(x_f - x_i)}{t_{\text{go}}^2}.
        \label{eq:apollo_k2}
\end{align}

Since the downrange target position is not known a~priori, it is estimated from
the current orbital elements by predicting where the unpowered trajectory would
intersect the target altitude.

\paragraph{Coefficient freezing.}
As $t_{\text{go}} \to 0$ the denominators $t_{\text{go}}^3$ and
$t_{\text{go}}^2$ drive the coefficients toward very large values.  To prevent
numerical instability the coefficients are frozen at their last computed values
once $t_{\text{go}}$ falls below a configurable threshold (nominally $10$\,s).

\paragraph{Steering-angle extraction.}
The total acceleration commanded by Eqs.~\eqref{eq:apollo_ax}--\eqref{eq:apollo_ay}
includes the gravitational contribution.  The thrust-only acceleration is
obtained by subtracting the local gravity vector, and the angle of attack is
then
\begin{equation}
  \alpha = \arctan\!\left(\frac{a_y^{\text{thrust}}}{a_x^{\text{thrust}}}\right)
           - \arctan\!\left(\frac{v_y}{v_x}\right).
  \label{eq:apollo_alpha}
\end{equation}

An optional thrust-magnitude-control mode allows the engine throttle to follow
the commanded acceleration magnitude
$|\mathbf{a}_{\text{thrust}}| = \sqrt{(a_x^{\text{thrust}})^2 +
(a_y^{\text{thrust}})^2}$, converting it to a thrust force via $F_T = m\,
|\mathbf{a}_{\text{thrust}}|$, limited by the maximum available thrust.

% -----------------------------------------------------------------------------
\subsection{Trajectory Optimisation}
\label{ssec:optimisation}

\subsubsection{Problem Formulation}
\label{sssec:opt_formulation}

The optimisation seeks the initial kick angle $\alpha_k$ that minimises the
total propellant consumption required to reach the target circular orbit:
\begin{equation}
  \min_{\alpha_k}\; J(\alpha_k)
    = m_{\text{prop,ascent}}(\alpha_k)
    + m_{\text{prop,circ}}(\alpha_k).
  \label{eq:obj}
\end{equation}

The single design variable is bounded:
$\alpha_k \in [\alpha_{\min},\,\alpha_{\max}]$ (nominally $[-5°,\,-2°]$).

Three soft constraints are enforced through a large penalty value
($J = 10^9$\,kg) applied when any condition is violated:
\begin{enumerate}
  \item propellant required for circularisation exceeds the available supply,
  \item the circularisation burn duration exceeds a maximum limit
        (nominally $15$\,s),
  \item the achieved apogee deviates from the target by more than $0.2\,\%$.
\end{enumerate}

\subsubsection{Optimisation Algorithm}
\label{sssec:opt_algorithm}

Because the objective function is non-smooth --- discrete events such as
staging, engine ignitions, and cutoffs introduce discontinuities --- a
gradient-free approach is used.  A brute-force grid search evaluates
$J(\alpha_k)$ at $N = 1000$ uniformly spaced points across
$[\alpha_{\min},\,\alpha_{\max}]$ and selects the global minimum.  No local
refinement step is applied.

The grid resolution is approximately $0.003°$, which is sufficient for
engineering accuracy.  With each trajectory evaluation taking on the order of
one second, the full optimisation completes in roughly $15$--$20$~minutes.

A key architectural feature is that the optimisation loop is entirely
independent of the guidance mode.  The outer loop supplies a candidate kick
angle and receives back a scalar cost; the inner simulation selects and
executes whichever guidance law is active.  Consequently the same search
procedure applies unchanged to all five guidance modes.

\subsubsection{Event-Based Engine Cutoff}
\label{sssec:seco_event}

The timing of SECO is not a free optimisation variable.  Instead, it is
determined implicitly during trajectory integration through an event function:
\begin{equation}
  g(t,\mathbf{x}) = r_{\text{apo}}(t,\mathbf{x}) - r_{\text{target}},
  \label{eq:seco_event}
\end{equation}
where $r_{\text{apo}}$ is the instantaneous apogee computed from the current
orbital elements (Eq.~\ref{eq:apsides}).  The integrator monitors $g$ during
the Stage~2 burn (above $65$\,km altitude); a sign change triggers root-finding
to locate the exact cutoff time.

The physical interpretation is direct: the engine burns until the trajectory,
if left unpowered from that instant onward, would coast to exactly the target
apogee altitude.  This event-driven approach automatically adapts the burn
duration to the trajectory shape produced by each guidance law and kick angle.

% -----------------------------------------------------------------------------
\subsection{Space Launcher Configuration}
\label{ssec:launcher}

The launch vehicle modelled in the simulator is representative of a
two-stage Falcon~9-class rocket.  Table~\ref{tab:launcher} summarises
the key parameters.

\begin{table}[htbp]
  \centering
  \caption{Launch-vehicle parameters used in the simulator.}
  \label{tab:launcher}
  \begin{tabular}{l c c c}
    \toprule
    \textbf{Parameter} & \textbf{Stage~1} & \textbf{Stage~2} & \textbf{Unit} \\
    \midrule
    Thrust (vacuum)       & 7\,600  & 934    & kN   \\
    Specific impulse      & 283     & 348    & s    \\
    Structural mass       & 25\,600 & 3\,900 & kg   \\
    Propellant mass       & 395\,700& 92\,670& kg   \\
    \bottomrule
  \end{tabular}
\end{table}

Stage separation occurs $3$\,s after MECO, and the second-stage engine ignites
$8$\,s after MECO, giving a $5$\,s coast interval between stages.

The aerodynamic reference area is $A = 10.52$\,m$^2$ with a constant drag
coefficient $C_D = 0.3$.

\subsubsection{Mission Parameters}

The baseline mission targets a $500$\,km circular orbit at an inclination of
$51.6°$ (representative of the International Space Station orbit), launching
from a latitude of $28.5°$ (Kennedy Space Center).  The resulting initial mass
at lift-off is approximately $517.9$\,t (sum of both stages, with zero payload
in the reference configuration).  Table~\ref{tab:mission} lists the nominal
mission parameters.

\begin{table}[htbp]
  \centering
  \caption{Baseline mission parameters.}
  \label{tab:mission}
  \begin{tabular}{l c c}
    \toprule
    \textbf{Parameter} & \textbf{Value} & \textbf{Unit} \\
    \midrule
    Target orbital altitude   & 500    & km   \\
    Target inclination        & 51.6   & deg  \\
    Launch latitude           & 28.5   & deg  \\
    Kick-angle search range   & $[-5,\,-2]$ & deg \\
    Guidance update rate      & 0.5    & s    \\
    Integration tolerance     & $10^{-8}$ & --- \\
    Maximum integration step  & 1.0    & s    \\
    Output time resolution    & 0.01   & s    \\
    \bottomrule
  \end{tabular}
\end{table}

% =============================================================================

\section{Simulator Implementation}
\label{sec:simulator_implementation}

This section describes how the trajectory simulation tool is structured and how
the theoretical principles introduced in the previous chapter are brought
together into a working ascent simulator.  The discussion covers the overall
software architecture, the reference frames adopted, the equations of motion
actually integrated, the environmental and force models, the guidance laws
available, and the optimisation strategy used to find the best initial
conditions.  Vehicle parameters representative of a Falcon~9-class launcher
close the section.

% -----------------------------------------------------------------------------
\subsection{Simulation Architecture and Workflow}
\label{ssec:sim_architecture}

The simulator is organised around three successive mission phases, each
resolved by a separate call to the numerical integrator:

\begin{enumerate}
  \item \textbf{Stage~1 powered ascent.}  From lift-off until first-stage
        propellant depletion (Main Engine Cut-Off, MECO).  During this phase
        the vehicle follows an initial vertical rise, a pitchover kick
        manoeuvre, and atmospheric flight with aerodynamic drag.

  \item \textbf{Stage~2 powered ascent with guidance.}  After stage separation
        and a brief coast interval, the second-stage engine ignites.  An active
        guidance law (selected from five available modes) shapes the trajectory
        until the event-triggered Second Engine Cut-Off (SECO), which occurs
        when the instantaneous apogee of the coasting trajectory matches the
        target orbital altitude.

  \item \textbf{Ballistic coast and circularisation.}  The vehicle coasts
        unpowered to apogee, where an impulsive circularisation burn is applied
        to achieve the target circular orbit.  A final short coast in stable
        orbit concludes the simulation.
\end{enumerate}

\subsubsection{State Vector}

The base state vector integrated by the solver is
\begin{equation}
  \mathbf{x} = \bigl[\,s,\; r,\; v,\; \gamma,\; m\,\bigr]^T,
  \label{eq:state_base}
\end{equation}
where $s$ is the downrange distance along the Earth surface, $r$ is the
geocentric radius, $v$ is the velocity magnitude in the rotating frame,
$\gamma$ is the flight-path angle, and $m$ is the vehicle mass.

When Earth-rotation effects are enabled, the state is extended to include
geocentric latitude~$\phi$:
\begin{equation}
  \mathbf{x} = \bigl[\,s,\; r,\; v,\; \gamma,\; m,\; \phi\,\bigr]^T.
  \label{eq:state_rot}
\end{equation}

If heading tracking is additionally activated, the local heading (azimuth)
$\psi$ becomes the seventh state component:
\begin{equation}
  \mathbf{x} = \bigl[\,s,\; r,\; v,\; \gamma,\; m,\; \phi,\; \psi\,\bigr]^T.
  \label{eq:state_heading}
\end{equation}

\subsubsection{Numerical Integration}

Each mission phase is propagated with an explicit Runge--Kutta~4(5) scheme
(Dormand--Prince variant) provided by \texttt{scipy.integrate.solve\_ivp}.  The
solver uses adaptive step-size control with an absolute tolerance of
$10^{-8}$, a maximum integration step of $1$~s, and a dense-output time
resolution of $0.01$~s.

Phase transitions are handled through \emph{event functions}: scalar functions
of the state that the integrator monitors for zero-crossings.  When a crossing
is detected the integrator performs root-finding to locate the exact transition
time and terminates the current phase.  Events used include propellant
depletion (MECO), stage separation timing, atmosphere exit, target-apogee
matching (SECO), and ground-collision safeguards.

\subsubsection{Architecture Overview}

Figure~\ref{fig:sim_architecture} shows the high-level data flow.  An outer
optimisation loop varies the initial kick angle supplied to the simulation
engine.  Inside each simulation call, the equations of motion are evaluated at
every integrator step using the current environmental models (gravity,
atmosphere) and the active guidance law to compute the steering command.  Event
functions run concurrently and trigger phase switches when their conditions are
met.  After the full trajectory is obtained, the optimiser evaluates the
objective function and proceeds to the next candidate.

\begin{figure}[htbp]
  \centering
  \begin{tikzpicture}[
      >=latex,
      block/.style  = {draw, rounded corners, minimum height=1.0cm,
                        minimum width=2.6cm, align=center, font=\small},
      smallblock/.style = {draw, rounded corners, minimum height=0.8cm,
                        minimum width=2.2cm, align=center, font=\footnotesize},
      arrow/.style  = {->, thick},
      darrow/.style = {<->, thick},
    ]

    % Optimisation loop
    \node[block, fill=blue!12] (opt) {Optimisation\\Loop};

    % Simulation engine
    \node[block, fill=green!10, right=2.2cm of opt] (sim) {Simulation\\Engine};

    % Inner modules (below sim)
    \node[smallblock, fill=orange!10, below left=1.4cm and -0.3cm of sim]
          (guid) {Guidance\\Module};
    \node[smallblock, fill=orange!10, below=1.4cm of sim]
          (eom) {Equations\\of Motion};
    \node[smallblock, fill=orange!10, below right=1.4cm and -0.3cm of sim]
          (env) {Environment\\Models};

    % Event detection
    \node[block, fill=red!10, right=2.2cm of sim] (evt) {Event\\Detection};

    % Outputs
    \node[block, fill=gray!12, right=2.2cm of evt] (out) {Trajectory\\Outputs};

    % Arrows
    \draw[arrow] (opt) -- node[above, font=\footnotesize]{$\alpha_{\text{kick}}$} (sim);
    \draw[arrow] (sim) -- (evt);
    \draw[arrow] (evt) -- (out);
    \draw[arrow] (out.south) -- ++(0,-1.2) -| node[near start, below, font=\footnotesize]{$J(\alpha)$} (opt.south);

    % Inner connections
    \draw[arrow] (sim) -- (guid);
    \draw[arrow] (sim) -- (eom);
    \draw[arrow] (sim) -- (env);
    \draw[arrow] (guid) -- (eom);
    \draw[arrow] (env) -- (eom);

  \end{tikzpicture}
  \caption{High-level architecture of the trajectory simulator.  The
           optimisation loop feeds candidate kick angles into the simulation
           engine, which dispatches environment evaluations and guidance
           commands to the equations of motion at each integration step.  Event
           detection triggers phase transitions and engine cutoff.  The
           resulting trajectory is scored and fed back to the optimiser.}
  \label{fig:sim_architecture}
\end{figure}

% -----------------------------------------------------------------------------
\subsection{Reference Frames}
\label{ssec:reference_frames}

The simulator employs a mixed-frame strategy that keeps guidance and
atmosphere coupling straightforward while preserving correct orbital-mechanics
metrics.

\begin{itemize}
  \item \textbf{Rotating frame (ECEF-like).}  All ascent dynamics quantities
        --- velocity magnitude~$v$, flight-path angle~$\gamma$, and heading
        decomposition --- are expressed in a frame that co-rotates with the
        Earth.  This simplifies the treatment of aerodynamic drag, which
        depends on the velocity relative to the atmosphere.

  \item \textbf{Inertial frame (ECI).}  Orbital-energy and orbital-element
        calculations (semi-major axis, eccentricity, apogee radius) are
        performed in an Earth-Centred Inertial frame.  This is essential for
        correct event detection (e.g.\ target-apogee matching) and end-of-run
        orbit characterisation.

  \item \textbf{Local East--North--Up (ENU).}  Pseudo-force accelerations
        (Coriolis and centrifugal) are first evaluated in the local ENU basis at
        the vehicle's current latitude, then projected onto the along-heading
        and radial directions used by the equations of motion.
\end{itemize}

\subsubsection{Frame Transition During Flight}

During powered ascent and early coasting the integration proceeds in the
rotating frame with optional pseudo-force corrections.  Once the vehicle
reaches the final ballistic coast phase (post-SECO), the propagation switches
to pure inertial dynamics: pseudo-force terms are deactivated to avoid
introducing spurious accelerations when the trajectory is governed solely by
two-body gravity.

\subsubsection{ECEF-to-ECI Velocity Conversion}

At any instant the local rotating-frame velocity components are converted to
their inertial counterparts by adding the Earth's rotational eastward speed.
The horizontal velocity magnitude is $v_h = v\cos\gamma$, which decomposes as
\begin{equation}
  v_N = v_h\cos\psi, \qquad
  v_E^{\text{ECEF}} = v_h\sin\psi.
  \label{eq:vel_ecef}
\end{equation}

Adding Earth rotation yields
\begin{equation}
  v_E^{\text{ECI}} = v_E^{\text{ECEF}} + \Omega_E\, r\cos\phi,
  \label{eq:ve_eci}
\end{equation}
from which the inertial speed and flight-path angle follow:
\begin{align}
  v_{\text{ECI}} &= \sqrt{\bigl(v_E^{\text{ECI}}\bigr)^2 + v_N^2
                    + \bigl(v\sin\gamma\bigr)^2}, \label{eq:v_eci} \\[4pt]
  \gamma_{\text{ECI}} &= \operatorname{atan2}\!\bigl(v\sin\gamma,\;
                    \sqrt{\bigl(v_E^{\text{ECI}}\bigr)^2 + v_N^2}\;\bigr).
                    \label{eq:gamma_eci}
\end{align}

These inertial quantities are used whenever orbital elements are computed.

% -----------------------------------------------------------------------------
\subsection{Launch Azimuth and Heading}
\label{ssec:azimuth_heading}

\subsubsection{Launch Azimuth Determination}
\label{sssec:launch_azimuth}

The launch azimuth defines the initial heading of the ascent trajectory in the
horizontal plane and directly controls the achievable orbital inclination.
Two modes are available.

\paragraph{Geometric (inertial) azimuth.}
Given a target inclination~$i$ and launch latitude~$\phi$, the classical
spherical-geometry relation gives the inertial azimuth measured from true
north:
\begin{equation}
  \sin\beta_{\text{geom}} = \frac{\cos i}{\cos\phi}.
  \label{eq:az_geom}
\end{equation}

\paragraph{Corrected (rotating-frame) azimuth.}
Because the rotating-frame velocity seen by the guidance system differs from
the inertial velocity by the eastward surface speed, the launch azimuth is
corrected:
\begin{equation}
  v_{\text{rot}} = \Omega_E\, R_E\cos\phi,
  \label{eq:vrot}
\end{equation}
\begin{align}
  v_N &= v_{\text{orb}}\cos\beta_{\text{geom}}, \label{eq:vn_corr} \\
  v_E^{\text{ECEF}} &= v_{\text{orb}}\sin\beta_{\text{geom}} - v_{\text{rot}},
  \label{eq:ve_corr}
\end{align}
where $v_{\text{orb}} = \sqrt{\mu/r_{\text{target}}}$ is the circular speed at
the target altitude.  The corrected azimuth is then
\begin{equation}
  \beta_{\text{corr}} = \operatorname{atan2}\!\bigl(v_E^{\text{ECEF}},\, v_N\bigr).
  \label{eq:az_corr}
\end{equation}

The simulator defaults to the corrected mode; a configuration flag may select
the geometric mode instead.

\subsubsection{Heading Propagation}
\label{sssec:heading_propagation}

When heading tracking is enabled, the local azimuth $\psi$ is propagated as a
state variable.  Its rate is computed from spherical-geometry consistency with
the great-circle ground track:
\begin{equation}
  \dot\psi = \tan\phi\;\tan\psi\;\dot\phi,
  \label{eq:psi_dot}
\end{equation}
where $\dot\phi$ is obtained from the latitude kinematic relation detailed in
Section~\ref{ssec:eom}.  An optional cross-heading pseudo-force contribution
may be added:
\begin{equation}
  \dot\psi \mathrel{+}= \frac{a_{\perp}}{v\cos\gamma},
  \label{eq:psi_pseudo}
\end{equation}
with $a_{\perp}$ the component of the Coriolis and centrifugal acceleration
perpendicular to the heading direction (see Section~\ref{sssec:pseudo_forces}).

\subsubsection{Achieved Inclination}
\label{sssec:achieved_inclination}

At the end of the simulation the achieved orbital inclination is evaluated by
converting the final rotating-frame velocity to inertial components using
Eqs.~\eqref{eq:vel_ecef}--\eqref{eq:ve_eci} and computing the inertial
velocity azimuth:
\begin{equation}
  \beta_{\text{ECI}} = \operatorname{atan2}\!\bigl(v_E^{\text{ECI}},\, v_N\bigr).
  \label{eq:beta_eci}
\end{equation}

The inclination follows from the classical relation:
\begin{equation}
  \cos i_{\text{achieved}} = \cos\phi\;\sin\beta_{\text{ECI}}.
  \label{eq:inc_achieved}
\end{equation}

The inclination drift $\Delta i = i_{\text{achieved}} - i_{\text{target}}$ is
a valuable metric for assessing the fidelity of the rotating-frame treatment
and the effectiveness of the heading-tracking augmentation.

% -----------------------------------------------------------------------------
\subsection{Equations of Motion}
\label{ssec:eom}

\subsubsection{Base Dynamics (Non-Rotating Form)}

The core differential system, expressed in spherical coordinates before
optional pseudo-force augmentation, is
\begin{align}
  \dot s      &= \frac{R_E}{r}\,v\cos\gamma,
               \label{eq:sdot}\\[3pt]
  \dot r      &= v\sin\gamma,
               \label{eq:rdot}\\[3pt]
  \dot v      &= \frac{F_T}{m}\cos\alpha
                 - \frac{F_D}{m}
                 - g(r)\sin\gamma,
               \label{eq:vdot}\\[3pt]
  \dot\gamma  &= \frac{1}{v}\!\left(
                   \frac{F_T}{m}\sin\alpha
                   + \frac{F_L}{m}
                   - \Bigl[g(r) - \frac{v^2}{r}\Bigr]\cos\gamma
                 \right),
               \label{eq:gammadot}\\[3pt]
  \dot m      &= -\frac{F_T}{I_{sp}\,g_0},
               \label{eq:mdot}
\end{align}
where $\alpha$ is the angle of attack (steering command from the guidance law),
$F_T$ and $I_{sp}$ are the current thrust and specific impulse, $F_D$ is
aerodynamic drag, $F_L$ is lift (set to zero in the present implementation),
and $g(r) = \mu/r^2$ is the gravitational acceleration.

The term $v^2/r$ in Eq.~\eqref{eq:gammadot} represents the centripetal
acceleration arising from the curved trajectory; when $v^2/r$ exceeds
$g\cos\gamma$ the net effect tends to raise the flight-path angle.

Two numerical safeguards are applied:
\begin{itemize}
  \item \emph{Small-velocity guard:} when $v < 10^{-6}$\,m/s the angular rate
        $\dot\gamma$ is forced to zero to prevent a singularity.
  \item \emph{Vertical-hold constraint:} during the initial vertical rise
        (before the pitchover kick begins) $\dot\gamma$ is explicitly set to
        zero regardless of geometric terms.
\end{itemize}

\subsubsection{Rotating-Frame Augmentation}

When Earth-rotation effects are enabled, Coriolis and centrifugal
accelerations computed in the local ENU frame (see
Section~\ref{sssec:pseudo_forces}) are projected onto the along-heading and
radial directions:
\begin{equation}
  a_{h,\parallel} = a_E\sin\psi + a_N\cos\psi,
  \label{eq:a_parallel}
\end{equation}
where $a_E$ and $a_N$ collect both Coriolis and centrifugal components.
The additive corrections to the velocity and flight-path-angle rates are
\begin{align}
  \Delta\dot v     &= a_{h,\parallel}\cos\gamma + a_{\text{up}}\sin\gamma,
                    \label{eq:dv_pseudo}\\[3pt]
  \Delta\dot\gamma &= \frac{-a_{h,\parallel}\sin\gamma
                      + a_{\text{up}}\cos\gamma}{v},
                    \label{eq:dgamma_pseudo}
\end{align}
so that the integrated rates become
$\dot v_{\text{tot}} = \dot v + \Delta\dot v$ and
$\dot\gamma_{\text{tot}} = \dot\gamma + \Delta\dot\gamma$.

\subsubsection{Latitude Propagation}

Rather than integrating a differential latitude rate directly, the simulator
reconstructs latitude from the downrange angle using great-circle geometry,
which avoids numerical drift:
\begin{equation}
  \sigma = \frac{s}{R_E}, \qquad
  \sin\phi = \sin\phi_0\cos\sigma + \cos\phi_0\sin\sigma\cos\beta_0,
  \label{eq:lat_recon}
\end{equation}
where $\phi_0$ is the launch latitude and $\beta_0$ is the inertial launch
azimuth.  The latitude rate required by the heading equation
\eqref{eq:psi_dot} is obtained by differentiating analytically:
\begin{equation}
  \dot\phi = \frac{-\sin\phi_0\sin\sigma
    + \cos\phi_0\cos\sigma\cos\beta_0}{\cos\phi}\;\frac{\dot s}{R_E}.
  \label{eq:phidot}
\end{equation}

% -----------------------------------------------------------------------------
\subsection{Environmental Models}
\label{ssec:env_models}

\subsubsection{Gravitational Model}
\label{sssec:gravity}

Gravity is modelled as a central inverse-square field:
\begin{equation}
  g(r) = \frac{\mu}{r^2},
  \label{eq:gravity}
\end{equation}
with $\mu = 3.986\times10^{14}$\,m$^3$\,s$^{-2}$.  The Earth is assumed
spherical; higher-order harmonics such as $J_2$ are not included.

\subsubsection{Atmospheric Model}
\label{sssec:atmosphere}

Atmospheric density follows an exponential profile:
\begin{equation}
  \rho(h) = \rho_0\,\exp\!\left(-\frac{h}{H}\right),
  \qquad h = r - R_E,
  \label{eq:atmosphere}
\end{equation}
with sea-level density $\rho_0 = 1.225$\,kg\,m$^{-3}$ and scale height
$H = 8500$\,m.  The dynamic pressure is $q = \tfrac{1}{2}\rho\,v^2$.

Two criteria are available for declaring atmosphere exit and
activating exo-atmospheric guidance:
\begin{itemize}
  \item \emph{Altitude threshold:} $h > 65$\,km.
  \item \emph{Dynamic-pressure threshold:} $q < 1000$\,Pa (configurable).
\end{itemize}
The selected criterion is a simulation parameter.

% -----------------------------------------------------------------------------
\subsection{External Forces}
\label{ssec:forces}

\subsubsection{Aerodynamic Forces}
\label{sssec:aero}

The aerodynamic drag force opposing the velocity vector is
\begin{equation}
  F_D = q\,C_D\,A = \tfrac{1}{2}\,\rho\,v^2\,C_D\,A,
  \label{eq:drag}
\end{equation}
with constant drag coefficient $C_D = 0.3$ and reference area $A = 10.52$\,m$^2$.
A lift model is available in the formulation (Eq.~\ref{eq:gammadot}), but in
all analyses presented here lift is neglected ($F_L = 0$).

\subsubsection{Propulsive Forces}
\label{sssec:propulsion}

Thrust magnitude $F_T$ and specific impulse $I_{sp}$ are stage-dependent
constants that change at discrete events (MECO, stage separation, second-stage
ignition, SECO).  Thrust is directed at an angle of attack $\alpha$ relative
to the velocity vector; $\alpha$ is the primary control variable and is
determined at each integration step by the active guidance law
(Section~\ref{ssec:guidance}).

The resulting vehicle mass depletion obeys
\begin{equation}
  \dot m = -\frac{F_T}{I_{sp}\,g_0},
  \label{eq:mdot_prop}
\end{equation}
where $g_0 = 9.80665$\,m\,s$^{-2}$ is the standard gravitational acceleration.

\subsubsection{Pseudo-Forces (Earth Rotation)}
\label{sssec:pseudo_forces}

When the integration is performed in the rotating frame with Earth-rotation
corrections enabled, inertial pseudo-accelerations are applied.

\paragraph{Coriolis acceleration.}
In the local ENU basis at latitude $\phi$:
\begin{align}
  a_E^{\text{cor}} &= -2\Omega_E\cos\phi\;v_{\text{up}}
                      + 2\Omega_E\sin\phi\;v_N,
                    \label{eq:cor_e} \\
  a_N^{\text{cor}} &= -2\Omega_E\sin\phi\;v_E,
                    \label{eq:cor_n} \\
  a_{\text{up}}^{\text{cor}} &= 2\Omega_E\cos\phi\;v_E.
                    \label{eq:cor_u}
\end{align}

\paragraph{Centrifugal acceleration.}
\begin{align}
  a_E^{\text{cen}} &= 0,
                    \label{eq:cen_e} \\
  a_N^{\text{cen}} &= -\Omega_E^2\,r\,\sin\phi\cos\phi,
                    \label{eq:cen_n} \\
  a_{\text{up}}^{\text{cen}} &= \Omega_E^2\,r\,\cos^2\!\phi.
                    \label{eq:cen_u}
\end{align}

\paragraph{Projection onto flight-dynamics directions.}
The total pseudo-acceleration in each ENU direction is summed ($a_E = a_E^{\text{cor}} + a_E^{\text{cen}}$, etc.) and then projected onto the heading direction and the vertical:
\begin{align}
  a_{h,\parallel} &= a_E\sin\psi + a_N\cos\psi
    & &\text{(along-heading)}, \label{eq:a_par} \\
  a_{h,\perp}     &= a_E\cos\psi - a_N\sin\psi
    & &\text{(cross-heading)}, \label{eq:a_perp}
\end{align}
with $a_{\text{up}}$ retained for the radial direction.

The along-heading and up components enter the velocity and flight-path-angle
corrections (Eqs.~\ref{eq:dv_pseudo}--\ref{eq:dgamma_pseudo}).  If cross-heading
pseudo-force tracking is enabled, $a_{h,\perp}$ contributes to the heading
rate via Eq.~\eqref{eq:psi_pseudo}.

% -----------------------------------------------------------------------------
\subsection{Performance Metrics}
\label{ssec:performance}

\subsubsection{Orbital Elements}
\label{sssec:orbital_elements}

Once the velocity is converted to the inertial frame
(Eqs.~\ref{eq:v_eci}--\ref{eq:gamma_eci}), the standard two-body orbital
elements are computed:
\begin{align}
  a &= \frac{\mu\,r}{2\mu - r\,v_{\text{ECI}}^2},
    \label{eq:sma}\\[3pt]
  e &= \sqrt{1 - \frac{\bigl(r\,v_{\text{ECI}}\cos\gamma_{\text{ECI}}\bigr)^2}
       {\mu\,a}},
    \label{eq:ecc}\\[3pt]
  r_{\text{apo}} &= a\,(1+e), \qquad
  r_{\text{peri}} = a\,(1-e).
  \label{eq:apsides}
\end{align}

These quantities are evaluated continuously during Stage~2 powered flight: the
apogee radius $r_{\text{apo}}$ is compared with the target orbital radius in
the SECO event function (Section~\ref{sssec:seco_event}).

\subsubsection{Circularisation \texorpdfstring{$\Delta v$}{Delta-v}}
\label{sssec:circ_dv}

At apogee the velocity deficit with respect to a circular orbit is
\begin{equation}
  \Delta v_{\text{circ}} = \left|v_{\text{circ}} - v_{\text{apo}}\right|,
  \qquad v_{\text{circ}} = \sqrt{\frac{\mu}{r_{\text{target}}}},
  \label{eq:dv_circ}
\end{equation}
where $v_{\text{apo}}$ is the coast-phase velocity at apogee.  The propellant
mass required for this impulsive burn follows from the Tsiolkovsky equation:
\begin{equation}
  m_{\text{prop,circ}} = m_{\text{apo}}\!\left(1 - \exp\!\left[
    -\frac{\Delta v_{\text{circ}}}{I_{sp}\,g_0}\right]\right).
  \label{eq:mprop_circ}
\end{equation}

\subsubsection{Trajectory Losses and Gains}
\label{sssec:losses_gains}

The total velocity change delivered by the propulsion system exceeds the ideal
circular-orbit speed because several loss mechanisms are present during the
ascent:

\begin{itemize}
  \item \textbf{Gravity losses} arise because a component of gravity opposes the
        velocity during non-horizontal flight, so that not all thrust translates
        into kinetic-energy gain. These losses are most significant during the
        steep early ascent and are on the order of
        $1000$\,m\,s$^{-1}$ for a typical mission.

  \item \textbf{Drag losses} result from aerodynamic deceleration in the
        lower atmosphere, typically contributing ${\sim}100$\,m\,s$^{-1}$.

  \item \textbf{Steering losses} occur when the thrust vector is not aligned
        with the velocity direction ($\alpha \neq 0$), reducing the effective
        thrust along the trajectory.  Their magnitude depends on the guidance
        law and can reach several hundred m\,s$^{-1}$.
\end{itemize}

A positive contribution comes from \textbf{Earth rotation}: the launch site's
eastward surface velocity provides a free initial velocity increment
\begin{equation}
  \Delta v_{\text{rot}} = v_{\text{rot}}\sin\beta
  = \Omega_E\,R_E\cos\phi\;\sin\beta,
  \label{eq:dv_rot}
\end{equation}
where $\beta$ is the launch azimuth.  For an eastward launch from latitude
$28.5°$ this gain is approximately $410$\,m\,s$^{-1}$, a non-negligible
fraction of the orbital velocity budget.

% -----------------------------------------------------------------------------
\subsection{Guidance Laws}
\label{ssec:guidance}

Five guidance strategies are available in the simulator.  All share a common
initial pitchover manoeuvre; they differ in how the angle of attack~$\alpha$ is
computed after the vehicle exits the atmosphere.

\subsubsection{Common Initial Kick Phase}
\label{sssec:kick_phase}

Regardless of the guidance mode selected, the first seconds of flight follow
the same pitch programme.  The vehicle rises vertically until time $T_k$
(nominally $7.5$\,s), then executes a symmetric triangular pitchover manoeuvre
over a duration $\Delta T_k$ ($45$\,s):
\begin{equation}
  \alpha(t) =
  \begin{cases}
    0, & t < T_k, \\[2pt]
    \displaystyle\alpha_k\;\frac{2(t - T_k)}{\Delta T_k},
      & T_k \le t < T_k + \tfrac{\Delta T_k}{2}, \\[6pt]
    \displaystyle\alpha_k\;2\!\left(1 - \frac{t - T_k}{\Delta T_k}\right),
      & T_k + \tfrac{\Delta T_k}{2} \le t < T_k + \Delta T_k, \\[6pt]
    0, & t \ge T_k + \Delta T_k,
  \end{cases}
  \label{eq:kick_profile}
\end{equation}
where $\alpha_k$ is the \emph{initial kick angle}, the single design variable
optimised by the trajectory optimiser (Section~\ref{ssec:optimisation}).
Typical optimum values lie in the range $\alpha_k \in [-5°,\,-2°]$.

After the kick concludes and until the atmosphere-exit criterion is met,
$\alpha$ is held at zero.  Active guidance is then engaged.

\subsubsection{Gravity Turn}
\label{sssec:gt}

The gravity-turn mode applies no active steering after the initial kick:
\begin{equation}
  \alpha(t) = 0, \qquad \forall\; t > T_k + \Delta T_k.
  \label{eq:gt}
\end{equation}

The trajectory is shaped exclusively by gravity and the initial conditions.
This constitutes the simplest guidance baseline.

\subsubsection{Simple Polynomial}
\label{sssec:simple_poly}

Upon atmosphere exit, a linear steering law drives the flight-path angle toward
zero (horizontal flight for circular-orbit insertion):
\begin{equation}
  \alpha(\tau) = a_0 + a_1\,\tau,
  \label{eq:simple_poly}
\end{equation}
where $\tau = t_{\text{go}} / 100$ is a normalised time-to-go parameter, and
the coefficients are determined from the boundary conditions
\begin{equation}
  a_0 = \gamma_{\text{target}} = 0, \qquad
  a_1 = \gamma_{\text{current}}.
  \label{eq:simple_poly_bc}
\end{equation}

The time-to-go is estimated from the remaining altitude and the current radial
velocity component.  Coefficients are recomputed at a guidance update rate of
$0.5$\,s.

\subsubsection{Linear Tangent Steering}
\label{sssec:lts}

Linear tangent steering prescribes that the tangent of the total velocity
inclination angle $(\alpha + \gamma)$ varies linearly with time-to-go:
\begin{equation}
  \tan\bigl(\alpha + \gamma\bigr) = a\;t_{\text{go}} + b.
  \label{eq:lts}
\end{equation}

The boundary conditions are:
\begin{itemize}
  \item Terminal horizontal flight: $\tan(\alpha_f + \gamma_f) = 0$
        at $t_{\text{go}} = 0$, giving $b = 0$.
  \item Current-state continuity: at the present $t_{\text{go}}$,
        $\tan(\gamma) = a\;t_{\text{go}}$, giving
        $a = \tan\gamma / t_{\text{go}}$.
\end{itemize}

The commanded angle of attack is therefore
\begin{equation}
  \alpha = \arctan\!\bigl(a\;t_{\text{go}}\bigr) - \gamma.
  \label{eq:lts_alpha}
\end{equation}

This classical method provides smooth, monotonic inclination control and
is updated every $0.5$\,s.

\subsubsection{Bilinear Tangent Steering}
\label{sssec:bts}

The bilinear tangent law generalises linear tangent steering by representing
the inclination-angle tangent as a ratio of two linear functions of
$t_{\text{go}}$:
\begin{equation}
  \tan\bigl(\alpha + \gamma\bigr)
    = \frac{c_1\,t_{\text{go}} + c_2}
           {c_1'\,t_{\text{go}} + c_2'}.
  \label{eq:bts}
\end{equation}

The four coefficients are determined from four boundary conditions:
\begin{enumerate}
  \item \emph{Terminal value}: $c_2 = 0$,
        enforcing $\tan(\alpha_f + \gamma_f) = 0$ at $t_{\text{go}} = 0$.
  \item \emph{Terminal rate}: $c_1$ is set to a small constant
        (nominally $0.02$) to ensure a smooth, non-abrupt approach to
        horizontal flight.
  \item \emph{Normalisation}: $c_2' = 1$ removes the scaling ambiguity.
  \item \emph{Initial value}: $c_1'$ is computed so that the inclination
        at the current $t_{\text{go}}$ matches a blended target:
        \begin{equation}
          \gamma_{\text{blended}} = 0.7\;\gamma_{\text{current}}
            + 0.3\;\gamma_{\text{targeting}},
          \label{eq:gamma_blend}
        \end{equation}
        where $\gamma_{\text{targeting}}$ is a geometric estimate of the
        flight-path angle required to reach the target altitude over the
        remaining downrange distance.
\end{enumerate}

The bilinear form controls both the value and the rate of change of the
inclination angle at the terminal point, providing smoother orbit insertion
than the purely linear variant.

\subsubsection{Apollo Polynomial Guidance}
\label{sssec:apollo}

The most sophisticated guidance mode follows the explicit polynomial approach
originating from the Apollo programme.  Two independent channels command the
total acceleration in the downrange ($x$) and altitude ($y$) directions:
\begin{align}
  a_x(t) &= k_1\,(t - t_{\text{epoch}}) + k_2, \label{eq:apollo_ax} \\
  a_y(t) &= k_3\,(t - t_{\text{epoch}}) + k_4, \label{eq:apollo_ay}
\end{align}
where $t_{\text{epoch}}$ is the time at which the coefficients were last
recomputed.

\paragraph{Terminal constraints.}
The four coefficients $k_1,\dots,k_4$ are chosen to satisfy position and
velocity constraints at the target orbital altitude:
\begin{itemize}
  \item reach the target geocentric radius $r_f$,
  \item attain the local circular-orbit velocity
        $v_f = \sqrt{\mu / r_f}$ in the horizontal direction,
  \item achieve zero radial velocity ($\gamma_f = 0$).
\end{itemize}

The resulting expressions, derived from polynomial interpolation over the
interval $t_{\text{go}}$, are (shown for one generic channel with initial
state subscript~$i$ and final subscript~$f$):
\begin{align}
  k_1 &= \frac{6\,(v_f + v_i)\,t_{\text{go}}
          - 12\,(x_f - x_i)}{t_{\text{go}}^3},
        \label{eq:apollo_k1} \\[4pt]
  k_2 &= \frac{-2\,(v_f + 2\,v_i)\,t_{\text{go}}
          + 6\,(x_f - x_i)}{t_{\text{go}}^2}.
        \label{eq:apollo_k2}
\end{align}

Since the downrange target position is not known a~priori, it is estimated from
the current orbital elements by predicting where the unpowered trajectory would
intersect the target altitude.

\paragraph{Coefficient freezing.}
As $t_{\text{go}} \to 0$ the denominators $t_{\text{go}}^3$ and
$t_{\text{go}}^2$ drive the coefficients toward very large values.  To prevent
numerical instability the coefficients are frozen at their last computed values
once $t_{\text{go}}$ falls below a configurable threshold (nominally $10$\,s).

\paragraph{Steering-angle extraction.}
The total acceleration commanded by Eqs.~\eqref{eq:apollo_ax}--\eqref{eq:apollo_ay}
includes the gravitational contribution.  The thrust-only acceleration is
obtained by subtracting the local gravity vector, and the angle of attack is
then
\begin{equation}
  \alpha = \arctan\!\left(\frac{a_y^{\text{thrust}}}{a_x^{\text{thrust}}}\right)
           - \arctan\!\left(\frac{v_y}{v_x}\right).
  \label{eq:apollo_alpha}
\end{equation}

An optional thrust-magnitude-control mode allows the engine throttle to follow
the commanded acceleration magnitude
$|\mathbf{a}_{\text{thrust}}| = \sqrt{(a_x^{\text{thrust}})^2 +
(a_y^{\text{thrust}})^2}$, converting it to a thrust force via $F_T = m\,
|\mathbf{a}_{\text{thrust}}|$, limited by the maximum available thrust.

% -----------------------------------------------------------------------------
\subsection{Trajectory Optimisation}
\label{ssec:optimisation}

\subsubsection{Problem Formulation}
\label{sssec:opt_formulation}

The optimisation seeks the initial kick angle $\alpha_k$ that minimises the
total propellant consumption required to reach the target circular orbit:
\begin{equation}
  \min_{\alpha_k}\; J(\alpha_k)
    = m_{\text{prop,ascent}}(\alpha_k)
    + m_{\text{prop,circ}}(\alpha_k).
  \label{eq:obj}
\end{equation}

The single design variable is bounded:
$\alpha_k \in [\alpha_{\min},\,\alpha_{\max}]$ (nominally $[-5°,\,-2°]$).

Three soft constraints are enforced through a large penalty value
($J = 10^9$\,kg) applied when any condition is violated:
\begin{enumerate}
  \item propellant required for circularisation exceeds the available supply,
  \item the circularisation burn duration exceeds a maximum limit
        (nominally $15$\,s),
  \item the achieved apogee deviates from the target by more than $0.2\,\%$.
\end{enumerate}

\subsubsection{Optimisation Algorithm}
\label{sssec:opt_algorithm}

Because the objective function is non-smooth --- discrete events such as
staging, engine ignitions, and cutoffs introduce discontinuities --- a
gradient-free approach is used.  A brute-force grid search evaluates
$J(\alpha_k)$ at $N = 1000$ uniformly spaced points across
$[\alpha_{\min},\,\alpha_{\max}]$ and selects the global minimum.  No local
refinement step is applied.

The grid resolution is approximately $0.003°$, which is sufficient for
engineering accuracy.  With each trajectory evaluation taking on the order of
one second, the full optimisation completes in roughly $15$--$20$~minutes.

A key architectural feature is that the optimisation loop is entirely
independent of the guidance mode.  The outer loop supplies a candidate kick
angle and receives back a scalar cost; the inner simulation selects and
executes whichever guidance law is active.  Consequently the same search
procedure applies unchanged to all five guidance modes.

\subsubsection{Event-Based Engine Cutoff}
\label{sssec:seco_event}

The timing of SECO is not a free optimisation variable.  Instead, it is
determined implicitly during trajectory integration through an event function:
\begin{equation}
  g(t,\mathbf{x}) = r_{\text{apo}}(t,\mathbf{x}) - r_{\text{target}},
  \label{eq:seco_event}
\end{equation}
where $r_{\text{apo}}$ is the instantaneous apogee computed from the current
orbital elements (Eq.~\ref{eq:apsides}).  The integrator monitors $g$ during
the Stage~2 burn (above $65$\,km altitude); a sign change triggers root-finding
to locate the exact cutoff time.

The physical interpretation is direct: the engine burns until the trajectory,
if left unpowered from that instant onward, would coast to exactly the target
apogee altitude.  This event-driven approach automatically adapts the burn
duration to the trajectory shape produced by each guidance law and kick angle.

% -----------------------------------------------------------------------------
\subsection{Space Launcher Configuration}
\label{ssec:launcher}

The launch vehicle modelled in the simulator is representative of a
two-stage Falcon~9-class rocket.  Table~\ref{tab:launcher} summarises
the key parameters.

\begin{table}[htbp]
  \centering
  \caption{Launch-vehicle parameters used in the simulator.}
  \label{tab:launcher}
  \begin{tabular}{l c c c}
    \toprule
    \textbf{Parameter} & \textbf{Stage~1} & \textbf{Stage~2} & \textbf{Unit} \\
    \midrule
    Thrust (vacuum)       & 7\,600  & 934    & kN   \\
    Specific impulse      & 283     & 348    & s    \\
    Structural mass       & 25\,600 & 3\,900 & kg   \\
    Propellant mass       & 395\,700& 92\,670& kg   \\
    \bottomrule
  \end{tabular}
\end{table}

Stage separation occurs $3$\,s after MECO, and the second-stage engine ignites
$8$\,s after MECO, giving a $5$\,s coast interval between stages.

The aerodynamic reference area is $A = 10.52$\,m$^2$ with a constant drag
coefficient $C_D = 0.3$.

\subsubsection{Mission Parameters}

The baseline mission targets a $500$\,km circular orbit at an inclination of
$51.6°$ (representative of the International Space Station orbit), launching
from a latitude of $28.5°$ (Kennedy Space Center).  The resulting initial mass
at lift-off is approximately $517.9$\,t (sum of both stages, with zero payload
in the reference configuration).  Table~\ref{tab:mission} lists the nominal
mission parameters.

\begin{table}[htbp]
  \centering
  \caption{Baseline mission parameters.}
  \label{tab:mission}
  \begin{tabular}{l c c}
    \toprule
    \textbf{Parameter} & \textbf{Value} & \textbf{Unit} \\
    \midrule
    Target orbital altitude   & 500    & km   \\
    Target inclination        & 51.6   & deg  \\
    Launch latitude           & 28.5   & deg  \\
    Kick-angle search range   & $[-5,\,-2]$ & deg \\
    Guidance update rate      & 0.5    & s    \\
    Integration tolerance     & $10^{-8}$ & --- \\
    Maximum integration step  & 1.0    & s    \\
    Output time resolution    & 0.01   & s    \\
    \bottomrule
  \end{tabular}
\end{table}

% =============================================================================

\section{Simulator Implementation}
\label{sec:simulator_implementation}

This section describes how the trajectory simulation tool is structured and how
the theoretical principles introduced in the previous chapter are brought
together into a working ascent simulator.  The discussion covers the overall
software architecture, the reference frames adopted, the equations of motion
actually integrated, the environmental and force models, the guidance laws
available, and the optimisation strategy used to find the best initial
conditions.  Vehicle parameters representative of a Falcon~9-class launcher
close the section.

% -----------------------------------------------------------------------------
\subsection{Simulation Architecture and Workflow}
\label{ssec:sim_architecture}

The simulator is organised around three successive mission phases, each
resolved by a separate call to the numerical integrator:

\begin{enumerate}
  \item \textbf{Stage~1 powered ascent.}  From lift-off until first-stage
        propellant depletion (Main Engine Cut-Off, MECO).  During this phase
        the vehicle follows an initial vertical rise, a pitchover kick
        manoeuvre, and atmospheric flight with aerodynamic drag.

  \item \textbf{Stage~2 powered ascent with guidance.}  After stage separation
        and a brief coast interval, the second-stage engine ignites.  An active
        guidance law (selected from five available modes) shapes the trajectory
        until the event-triggered Second Engine Cut-Off (SECO), which occurs
        when the instantaneous apogee of the coasting trajectory matches the
        target orbital altitude.

  \item \textbf{Ballistic coast and circularisation.}  The vehicle coasts
        unpowered to apogee, where an impulsive circularisation burn is applied
        to achieve the target circular orbit.  A final short coast in stable
        orbit concludes the simulation.
\end{enumerate}

\subsubsection{State Vector}

The base state vector integrated by the solver is
\begin{equation}
  \mathbf{x} = \bigl[\,s,\; r,\; v,\; \gamma,\; m\,\bigr]^T,
  \label{eq:state_base}
\end{equation}
where $s$ is the downrange distance along the Earth surface, $r$ is the
geocentric radius, $v$ is the velocity magnitude in the rotating frame,
$\gamma$ is the flight-path angle, and $m$ is the vehicle mass.

When Earth-rotation effects are enabled, the state is extended to include
geocentric latitude~$\phi$:
\begin{equation}
  \mathbf{x} = \bigl[\,s,\; r,\; v,\; \gamma,\; m,\; \phi\,\bigr]^T.
  \label{eq:state_rot}
\end{equation}

If heading tracking is additionally activated, the local heading (azimuth)
$\psi$ becomes the seventh state component:
\begin{equation}
  \mathbf{x} = \bigl[\,s,\; r,\; v,\; \gamma,\; m,\; \phi,\; \psi\,\bigr]^T.
  \label{eq:state_heading}
\end{equation}

\subsubsection{Numerical Integration}

Each mission phase is propagated with an explicit Runge--Kutta~4(5) scheme
(Dormand--Prince variant) provided by \texttt{scipy.integrate.solve\_ivp}.  The
solver uses adaptive step-size control with an absolute tolerance of
$10^{-8}$, a maximum integration step of $1$~s, and a dense-output time
resolution of $0.01$~s.

Phase transitions are handled through \emph{event functions}: scalar functions
of the state that the integrator monitors for zero-crossings.  When a crossing
is detected the integrator performs root-finding to locate the exact transition
time and terminates the current phase.  Events used include propellant
depletion (MECO), stage separation timing, atmosphere exit, target-apogee
matching (SECO), and ground-collision safeguards.

\subsubsection{Architecture Overview}

Figure~\ref{fig:sim_architecture} shows the high-level data flow.  An outer
optimisation loop varies the initial kick angle supplied to the simulation
engine.  Inside each simulation call, the equations of motion are evaluated at
every integrator step using the current environmental models (gravity,
atmosphere) and the active guidance law to compute the steering command.  Event
functions run concurrently and trigger phase switches when their conditions are
met.  After the full trajectory is obtained, the optimiser evaluates the
objective function and proceeds to the next candidate.

\begin{figure}[htbp]
  \centering
  \begin{tikzpicture}[
      >=latex,
      block/.style  = {draw, rounded corners, minimum height=1.0cm,
                        minimum width=2.6cm, align=center, font=\small},
      smallblock/.style = {draw, rounded corners, minimum height=0.8cm,
                        minimum width=2.2cm, align=center, font=\footnotesize},
      arrow/.style  = {->, thick},
      darrow/.style = {<->, thick},
    ]

    % Optimisation loop
    \node[block, fill=blue!12] (opt) {Optimisation\\Loop};

    % Simulation engine
    \node[block, fill=green!10, right=2.2cm of opt] (sim) {Simulation\\Engine};

    % Inner modules (below sim)
    \node[smallblock, fill=orange!10, below left=1.4cm and -0.3cm of sim]
          (guid) {Guidance\\Module};
    \node[smallblock, fill=orange!10, below=1.4cm of sim]
          (eom) {Equations\\of Motion};
    \node[smallblock, fill=orange!10, below right=1.4cm and -0.3cm of sim]
          (env) {Environment\\Models};

    % Event detection
    \node[block, fill=red!10, right=2.2cm of sim] (evt) {Event\\Detection};

    % Outputs
    \node[block, fill=gray!12, right=2.2cm of evt] (out) {Trajectory\\Outputs};

    % Arrows
    \draw[arrow] (opt) -- node[above, font=\footnotesize]{$\alpha_{\text{kick}}$} (sim);
    \draw[arrow] (sim) -- (evt);
    \draw[arrow] (evt) -- (out);
    \draw[arrow] (out.south) -- ++(0,-1.2) -| node[near start, below, font=\footnotesize]{$J(\alpha)$} (opt.south);

    % Inner connections
    \draw[arrow] (sim) -- (guid);
    \draw[arrow] (sim) -- (eom);
    \draw[arrow] (sim) -- (env);
    \draw[arrow] (guid) -- (eom);
    \draw[arrow] (env) -- (eom);

  \end{tikzpicture}
  \caption{High-level architecture of the trajectory simulator.  The
           optimisation loop feeds candidate kick angles into the simulation
           engine, which dispatches environment evaluations and guidance
           commands to the equations of motion at each integration step.  Event
           detection triggers phase transitions and engine cutoff.  The
           resulting trajectory is scored and fed back to the optimiser.}
  \label{fig:sim_architecture}
\end{figure}

% -----------------------------------------------------------------------------
\subsection{Reference Frames}
\label{ssec:reference_frames}

The simulator employs a mixed-frame strategy that keeps guidance and
atmosphere coupling straightforward while preserving correct orbital-mechanics
metrics.

\begin{itemize}
  \item \textbf{Rotating frame (ECEF-like).}  All ascent dynamics quantities
        --- velocity magnitude~$v$, flight-path angle~$\gamma$, and heading
        decomposition --- are expressed in a frame that co-rotates with the
        Earth.  This simplifies the treatment of aerodynamic drag, which
        depends on the velocity relative to the atmosphere.

  \item \textbf{Inertial frame (ECI).}  Orbital-energy and orbital-element
        calculations (semi-major axis, eccentricity, apogee radius) are
        performed in an Earth-Centred Inertial frame.  This is essential for
        correct event detection (e.g.\ target-apogee matching) and end-of-run
        orbit characterisation.

  \item \textbf{Local East--North--Up (ENU).}  Pseudo-force accelerations
        (Coriolis and centrifugal) are first evaluated in the local ENU basis at
        the vehicle's current latitude, then projected onto the along-heading
        and radial directions used by the equations of motion.
\end{itemize}

\subsubsection{Frame Transition During Flight}

During powered ascent and early coasting the integration proceeds in the
rotating frame with optional pseudo-force corrections.  Once the vehicle
reaches the final ballistic coast phase (post-SECO), the propagation switches
to pure inertial dynamics: pseudo-force terms are deactivated to avoid
introducing spurious accelerations when the trajectory is governed solely by
two-body gravity.

\subsubsection{ECEF-to-ECI Velocity Conversion}

At any instant the local rotating-frame velocity components are converted to
their inertial counterparts by adding the Earth's rotational eastward speed.
The horizontal velocity magnitude is $v_h = v\cos\gamma$, which decomposes as
\begin{equation}
  v_N = v_h\cos\psi, \qquad
  v_E^{\text{ECEF}} = v_h\sin\psi.
  \label{eq:vel_ecef}
\end{equation}

Adding Earth rotation yields
\begin{equation}
  v_E^{\text{ECI}} = v_E^{\text{ECEF}} + \Omega_E\, r\cos\phi,
  \label{eq:ve_eci}
\end{equation}
from which the inertial speed and flight-path angle follow:
\begin{align}
  v_{\text{ECI}} &= \sqrt{\bigl(v_E^{\text{ECI}}\bigr)^2 + v_N^2
                    + \bigl(v\sin\gamma\bigr)^2}, \label{eq:v_eci} \\[4pt]
  \gamma_{\text{ECI}} &= \operatorname{atan2}\!\bigl(v\sin\gamma,\;
                    \sqrt{\bigl(v_E^{\text{ECI}}\bigr)^2 + v_N^2}\;\bigr).
                    \label{eq:gamma_eci}
\end{align}

These inertial quantities are used whenever orbital elements are computed.

% -----------------------------------------------------------------------------
\subsection{Launch Azimuth and Heading}
\label{ssec:azimuth_heading}

\subsubsection{Launch Azimuth Determination}
\label{sssec:launch_azimuth}

The launch azimuth defines the initial heading of the ascent trajectory in the
horizontal plane and directly controls the achievable orbital inclination.
Two modes are available.

\paragraph{Geometric (inertial) azimuth.}
Given a target inclination~$i$ and launch latitude~$\phi$, the classical
spherical-geometry relation gives the inertial azimuth measured from true
north:
\begin{equation}
  \sin\beta_{\text{geom}} = \frac{\cos i}{\cos\phi}.
  \label{eq:az_geom}
\end{equation}

\paragraph{Corrected (rotating-frame) azimuth.}
Because the rotating-frame velocity seen by the guidance system differs from
the inertial velocity by the eastward surface speed, the launch azimuth is
corrected:
\begin{equation}
  v_{\text{rot}} = \Omega_E\, R_E\cos\phi,
  \label{eq:vrot}
\end{equation}
\begin{align}
  v_N &= v_{\text{orb}}\cos\beta_{\text{geom}}, \label{eq:vn_corr} \\
  v_E^{\text{ECEF}} &= v_{\text{orb}}\sin\beta_{\text{geom}} - v_{\text{rot}},
  \label{eq:ve_corr}
\end{align}
where $v_{\text{orb}} = \sqrt{\mu/r_{\text{target}}}$ is the circular speed at
the target altitude.  The corrected azimuth is then
\begin{equation}
  \beta_{\text{corr}} = \operatorname{atan2}\!\bigl(v_E^{\text{ECEF}},\, v_N\bigr).
  \label{eq:az_corr}
\end{equation}

The simulator defaults to the corrected mode; a configuration flag may select
the geometric mode instead.

\subsubsection{Heading Propagation}
\label{sssec:heading_propagation}

When heading tracking is enabled, the local azimuth $\psi$ is propagated as a
state variable.  Its rate is computed from spherical-geometry consistency with
the great-circle ground track:
\begin{equation}
  \dot\psi = \tan\phi\;\tan\psi\;\dot\phi,
  \label{eq:psi_dot}
\end{equation}
where $\dot\phi$ is obtained from the latitude kinematic relation detailed in
Section~\ref{ssec:eom}.  An optional cross-heading pseudo-force contribution
may be added:
\begin{equation}
  \dot\psi \mathrel{+}= \frac{a_{\perp}}{v\cos\gamma},
  \label{eq:psi_pseudo}
\end{equation}
with $a_{\perp}$ the component of the Coriolis and centrifugal acceleration
perpendicular to the heading direction (see Section~\ref{sssec:pseudo_forces}).

\subsubsection{Achieved Inclination}
\label{sssec:achieved_inclination}

At the end of the simulation the achieved orbital inclination is evaluated by
converting the final rotating-frame velocity to inertial components using
Eqs.~\eqref{eq:vel_ecef}--\eqref{eq:ve_eci} and computing the inertial
velocity azimuth:
\begin{equation}
  \beta_{\text{ECI}} = \operatorname{atan2}\!\bigl(v_E^{\text{ECI}},\, v_N\bigr).
  \label{eq:beta_eci}
\end{equation}

The inclination follows from the classical relation:
\begin{equation}
  \cos i_{\text{achieved}} = \cos\phi\;\sin\beta_{\text{ECI}}.
  \label{eq:inc_achieved}
\end{equation}

The inclination drift $\Delta i = i_{\text{achieved}} - i_{\text{target}}$ is
a valuable metric for assessing the fidelity of the rotating-frame treatment
and the effectiveness of the heading-tracking augmentation.

% -----------------------------------------------------------------------------
\subsection{Equations of Motion}
\label{ssec:eom}

\subsubsection{Base Dynamics (Non-Rotating Form)}

The core differential system, expressed in spherical coordinates before
optional pseudo-force augmentation, is
\begin{align}
  \dot s      &= \frac{R_E}{r}\,v\cos\gamma,
               \label{eq:sdot}\\[3pt]
  \dot r      &= v\sin\gamma,
               \label{eq:rdot}\\[3pt]
  \dot v      &= \frac{F_T}{m}\cos\alpha
                 - \frac{F_D}{m}
                 - g(r)\sin\gamma,
               \label{eq:vdot}\\[3pt]
  \dot\gamma  &= \frac{1}{v}\!\left(
                   \frac{F_T}{m}\sin\alpha
                   + \frac{F_L}{m}
                   - \Bigl[g(r) - \frac{v^2}{r}\Bigr]\cos\gamma
                 \right),
               \label{eq:gammadot}\\[3pt]
  \dot m      &= -\frac{F_T}{I_{sp}\,g_0},
               \label{eq:mdot}
\end{align}
where $\alpha$ is the angle of attack (steering command from the guidance law),
$F_T$ and $I_{sp}$ are the current thrust and specific impulse, $F_D$ is
aerodynamic drag, $F_L$ is lift (set to zero in the present implementation),
and $g(r) = \mu/r^2$ is the gravitational acceleration.

The term $v^2/r$ in Eq.~\eqref{eq:gammadot} represents the centripetal
acceleration arising from the curved trajectory; when $v^2/r$ exceeds
$g\cos\gamma$ the net effect tends to raise the flight-path angle.

Two numerical safeguards are applied:
\begin{itemize}
  \item \emph{Small-velocity guard:} when $v < 10^{-6}$\,m/s the angular rate
        $\dot\gamma$ is forced to zero to prevent a singularity.
  \item \emph{Vertical-hold constraint:} during the initial vertical rise
        (before the pitchover kick begins) $\dot\gamma$ is explicitly set to
        zero regardless of geometric terms.
\end{itemize}

\subsubsection{Rotating-Frame Augmentation}

When Earth-rotation effects are enabled, Coriolis and centrifugal
accelerations computed in the local ENU frame (see
Section~\ref{sssec:pseudo_forces}) are projected onto the along-heading and
radial directions:
\begin{equation}
  a_{h,\parallel} = a_E\sin\psi + a_N\cos\psi,
  \label{eq:a_parallel}
\end{equation}
where $a_E$ and $a_N$ collect both Coriolis and centrifugal components.
The additive corrections to the velocity and flight-path-angle rates are
\begin{align}
  \Delta\dot v     &= a_{h,\parallel}\cos\gamma + a_{\text{up}}\sin\gamma,
                    \label{eq:dv_pseudo}\\[3pt]
  \Delta\dot\gamma &= \frac{-a_{h,\parallel}\sin\gamma
                      + a_{\text{up}}\cos\gamma}{v},
                    \label{eq:dgamma_pseudo}
\end{align}
so that the integrated rates become
$\dot v_{\text{tot}} = \dot v + \Delta\dot v$ and
$\dot\gamma_{\text{tot}} = \dot\gamma + \Delta\dot\gamma$.

\subsubsection{Latitude Propagation}

Rather than integrating a differential latitude rate directly, the simulator
reconstructs latitude from the downrange angle using great-circle geometry,
which avoids numerical drift:
\begin{equation}
  \sigma = \frac{s}{R_E}, \qquad
  \sin\phi = \sin\phi_0\cos\sigma + \cos\phi_0\sin\sigma\cos\beta_0,
  \label{eq:lat_recon}
\end{equation}
where $\phi_0$ is the launch latitude and $\beta_0$ is the inertial launch
azimuth.  The latitude rate required by the heading equation
\eqref{eq:psi_dot} is obtained by differentiating analytically:
\begin{equation}
  \dot\phi = \frac{-\sin\phi_0\sin\sigma
    + \cos\phi_0\cos\sigma\cos\beta_0}{\cos\phi}\;\frac{\dot s}{R_E}.
  \label{eq:phidot}
\end{equation}

% -----------------------------------------------------------------------------
\subsection{Environmental Models}
\label{ssec:env_models}

\subsubsection{Gravitational Model}
\label{sssec:gravity}

Gravity is modelled as a central inverse-square field:
\begin{equation}
  g(r) = \frac{\mu}{r^2},
  \label{eq:gravity}
\end{equation}
with $\mu = 3.986\times10^{14}$\,m$^3$\,s$^{-2}$.  The Earth is assumed
spherical; higher-order harmonics such as $J_2$ are not included.

\subsubsection{Atmospheric Model}
\label{sssec:atmosphere}

Atmospheric density follows an exponential profile:
\begin{equation}
  \rho(h) = \rho_0\,\exp\!\left(-\frac{h}{H}\right),
  \qquad h = r - R_E,
  \label{eq:atmosphere}
\end{equation}
with sea-level density $\rho_0 = 1.225$\,kg\,m$^{-3}$ and scale height
$H = 8500$\,m.  The dynamic pressure is $q = \tfrac{1}{2}\rho\,v^2$.

Two criteria are available for declaring atmosphere exit and
activating exo-atmospheric guidance:
\begin{itemize}
  \item \emph{Altitude threshold:} $h > 65$\,km.
  \item \emph{Dynamic-pressure threshold:} $q < 1000$\,Pa (configurable).
\end{itemize}
The selected criterion is a simulation parameter.

% -----------------------------------------------------------------------------
\subsection{External Forces}
\label{ssec:forces}

\subsubsection{Aerodynamic Forces}
\label{sssec:aero}

The aerodynamic drag force opposing the velocity vector is
\begin{equation}
  F_D = q\,C_D\,A = \tfrac{1}{2}\,\rho\,v^2\,C_D\,A,
  \label{eq:drag}
\end{equation}
with constant drag coefficient $C_D = 0.3$ and reference area $A = 10.52$\,m$^2$.
A lift model is available in the formulation (Eq.~\ref{eq:gammadot}), but in
all analyses presented here lift is neglected ($F_L = 0$).

\subsubsection{Propulsive Forces}
\label{sssec:propulsion}

Thrust magnitude $F_T$ and specific impulse $I_{sp}$ are stage-dependent
constants that change at discrete events (MECO, stage separation, second-stage
ignition, SECO).  Thrust is directed at an angle of attack $\alpha$ relative
to the velocity vector; $\alpha$ is the primary control variable and is
determined at each integration step by the active guidance law
(Section~\ref{ssec:guidance}).

The resulting vehicle mass depletion obeys
\begin{equation}
  \dot m = -\frac{F_T}{I_{sp}\,g_0},
  \label{eq:mdot_prop}
\end{equation}
where $g_0 = 9.80665$\,m\,s$^{-2}$ is the standard gravitational acceleration.

\subsubsection{Pseudo-Forces (Earth Rotation)}
\label{sssec:pseudo_forces}

When the integration is performed in the rotating frame with Earth-rotation
corrections enabled, inertial pseudo-accelerations are applied.

\paragraph{Coriolis acceleration.}
In the local ENU basis at latitude $\phi$:
\begin{align}
  a_E^{\text{cor}} &= -2\Omega_E\cos\phi\;v_{\text{up}}
                      + 2\Omega_E\sin\phi\;v_N,
                    \label{eq:cor_e} \\
  a_N^{\text{cor}} &= -2\Omega_E\sin\phi\;v_E,
                    \label{eq:cor_n} \\
  a_{\text{up}}^{\text{cor}} &= 2\Omega_E\cos\phi\;v_E.
                    \label{eq:cor_u}
\end{align}

\paragraph{Centrifugal acceleration.}
\begin{align}
  a_E^{\text{cen}} &= 0,
                    \label{eq:cen_e} \\
  a_N^{\text{cen}} &= -\Omega_E^2\,r\,\sin\phi\cos\phi,
                    \label{eq:cen_n} \\
  a_{\text{up}}^{\text{cen}} &= \Omega_E^2\,r\,\cos^2\!\phi.
                    \label{eq:cen_u}
\end{align}

\paragraph{Projection onto flight-dynamics directions.}
The total pseudo-acceleration in each ENU direction is summed ($a_E = a_E^{\text{cor}} + a_E^{\text{cen}}$, etc.) and then projected onto the heading direction and the vertical:
\begin{align}
  a_{h,\parallel} &= a_E\sin\psi + a_N\cos\psi
    & &\text{(along-heading)}, \label{eq:a_par} \\
  a_{h,\perp}     &= a_E\cos\psi - a_N\sin\psi
    & &\text{(cross-heading)}, \label{eq:a_perp}
\end{align}
with $a_{\text{up}}$ retained for the radial direction.

The along-heading and up components enter the velocity and flight-path-angle
corrections (Eqs.~\ref{eq:dv_pseudo}--\ref{eq:dgamma_pseudo}).  If cross-heading
pseudo-force tracking is enabled, $a_{h,\perp}$ contributes to the heading
rate via Eq.~\eqref{eq:psi_pseudo}.

% -----------------------------------------------------------------------------
\subsection{Performance Metrics}
\label{ssec:performance}

\subsubsection{Orbital Elements}
\label{sssec:orbital_elements}

Once the velocity is converted to the inertial frame
(Eqs.~\ref{eq:v_eci}--\ref{eq:gamma_eci}), the standard two-body orbital
elements are computed:
\begin{align}
  a &= \frac{\mu\,r}{2\mu - r\,v_{\text{ECI}}^2},
    \label{eq:sma}\\[3pt]
  e &= \sqrt{1 - \frac{\bigl(r\,v_{\text{ECI}}\cos\gamma_{\text{ECI}}\bigr)^2}
       {\mu\,a}},
    \label{eq:ecc}\\[3pt]
  r_{\text{apo}} &= a\,(1+e), \qquad
  r_{\text{peri}} = a\,(1-e).
  \label{eq:apsides}
\end{align}

These quantities are evaluated continuously during Stage~2 powered flight: the
apogee radius $r_{\text{apo}}$ is compared with the target orbital radius in
the SECO event function (Section~\ref{sssec:seco_event}).

\subsubsection{Circularisation \texorpdfstring{$\Delta v$}{Delta-v}}
\label{sssec:circ_dv}

At apogee the velocity deficit with respect to a circular orbit is
\begin{equation}
  \Delta v_{\text{circ}} = \left|v_{\text{circ}} - v_{\text{apo}}\right|,
  \qquad v_{\text{circ}} = \sqrt{\frac{\mu}{r_{\text{target}}}},
  \label{eq:dv_circ}
\end{equation}
where $v_{\text{apo}}$ is the coast-phase velocity at apogee.  The propellant
mass required for this impulsive burn follows from the Tsiolkovsky equation:
\begin{equation}
  m_{\text{prop,circ}} = m_{\text{apo}}\!\left(1 - \exp\!\left[
    -\frac{\Delta v_{\text{circ}}}{I_{sp}\,g_0}\right]\right).
  \label{eq:mprop_circ}
\end{equation}

\subsubsection{Trajectory Losses and Gains}
\label{sssec:losses_gains}

The total velocity change delivered by the propulsion system exceeds the ideal
circular-orbit speed because several loss mechanisms are present during the
ascent:

\begin{itemize}
  \item \textbf{Gravity losses} arise because a component of gravity opposes the
        velocity during non-horizontal flight, so that not all thrust translates
        into kinetic-energy gain. These losses are most significant during the
        steep early ascent and are on the order of
        $1000$\,m\,s$^{-1}$ for a typical mission.

  \item \textbf{Drag losses} result from aerodynamic deceleration in the
        lower atmosphere, typically contributing ${\sim}100$\,m\,s$^{-1}$.

  \item \textbf{Steering losses} occur when the thrust vector is not aligned
        with the velocity direction ($\alpha \neq 0$), reducing the effective
        thrust along the trajectory.  Their magnitude depends on the guidance
        law and can reach several hundred m\,s$^{-1}$.
\end{itemize}

A positive contribution comes from \textbf{Earth rotation}: the launch site's
eastward surface velocity provides a free initial velocity increment
\begin{equation}
  \Delta v_{\text{rot}} = v_{\text{rot}}\sin\beta
  = \Omega_E\,R_E\cos\phi\;\sin\beta,
  \label{eq:dv_rot}
\end{equation}
where $\beta$ is the launch azimuth.  For an eastward launch from latitude
$28.5°$ this gain is approximately $410$\,m\,s$^{-1}$, a non-negligible
fraction of the orbital velocity budget.

% -----------------------------------------------------------------------------
\subsection{Guidance Laws}
\label{ssec:guidance}

Five guidance strategies are available in the simulator.  All share a common
initial pitchover manoeuvre; they differ in how the angle of attack~$\alpha$ is
computed after the vehicle exits the atmosphere.

\subsubsection{Common Initial Kick Phase}
\label{sssec:kick_phase}

Regardless of the guidance mode selected, the first seconds of flight follow
the same pitch programme.  The vehicle rises vertically until time $T_k$
(nominally $7.5$\,s), then executes a symmetric triangular pitchover manoeuvre
over a duration $\Delta T_k$ ($45$\,s):
\begin{equation}
  \alpha(t) =
  \begin{cases}
    0, & t < T_k, \\[2pt]
    \displaystyle\alpha_k\;\frac{2(t - T_k)}{\Delta T_k},
      & T_k \le t < T_k + \tfrac{\Delta T_k}{2}, \\[6pt]
    \displaystyle\alpha_k\;2\!\left(1 - \frac{t - T_k}{\Delta T_k}\right),
      & T_k + \tfrac{\Delta T_k}{2} \le t < T_k + \Delta T_k, \\[6pt]
    0, & t \ge T_k + \Delta T_k,
  \end{cases}
  \label{eq:kick_profile}
\end{equation}
where $\alpha_k$ is the \emph{initial kick angle}, the single design variable
optimised by the trajectory optimiser (Section~\ref{ssec:optimisation}).
Typical optimum values lie in the range $\alpha_k \in [-5°,\,-2°]$.

After the kick concludes and until the atmosphere-exit criterion is met,
$\alpha$ is held at zero.  Active guidance is then engaged.

\subsubsection{Gravity Turn}
\label{sssec:gt}

The gravity-turn mode applies no active steering after the initial kick:
\begin{equation}
  \alpha(t) = 0, \qquad \forall\; t > T_k + \Delta T_k.
  \label{eq:gt}
\end{equation}

The trajectory is shaped exclusively by gravity and the initial conditions.
This constitutes the simplest guidance baseline.

\subsubsection{Simple Polynomial}
\label{sssec:simple_poly}

Upon atmosphere exit, a linear steering law drives the flight-path angle toward
zero (horizontal flight for circular-orbit insertion):
\begin{equation}
  \alpha(\tau) = a_0 + a_1\,\tau,
  \label{eq:simple_poly}
\end{equation}
where $\tau = t_{\text{go}} / 100$ is a normalised time-to-go parameter, and
the coefficients are determined from the boundary conditions
\begin{equation}
  a_0 = \gamma_{\text{target}} = 0, \qquad
  a_1 = \gamma_{\text{current}}.
  \label{eq:simple_poly_bc}
\end{equation}

The time-to-go is estimated from the remaining altitude and the current radial
velocity component.  Coefficients are recomputed at a guidance update rate of
$0.5$\,s.

\subsubsection{Linear Tangent Steering}
\label{sssec:lts}

Linear tangent steering prescribes that the tangent of the total velocity
inclination angle $(\alpha + \gamma)$ varies linearly with time-to-go:
\begin{equation}
  \tan\bigl(\alpha + \gamma\bigr) = a\;t_{\text{go}} + b.
  \label{eq:lts}
\end{equation}

The boundary conditions are:
\begin{itemize}
  \item Terminal horizontal flight: $\tan(\alpha_f + \gamma_f) = 0$
        at $t_{\text{go}} = 0$, giving $b = 0$.
  \item Current-state continuity: at the present $t_{\text{go}}$,
        $\tan(\gamma) = a\;t_{\text{go}}$, giving
        $a = \tan\gamma / t_{\text{go}}$.
\end{itemize}

The commanded angle of attack is therefore
\begin{equation}
  \alpha = \arctan\!\bigl(a\;t_{\text{go}}\bigr) - \gamma.
  \label{eq:lts_alpha}
\end{equation}

This classical method provides smooth, monotonic inclination control and
is updated every $0.5$\,s.

\subsubsection{Bilinear Tangent Steering}
\label{sssec:bts}

The bilinear tangent law generalises linear tangent steering by representing
the inclination-angle tangent as a ratio of two linear functions of
$t_{\text{go}}$:
\begin{equation}
  \tan\bigl(\alpha + \gamma\bigr)
    = \frac{c_1\,t_{\text{go}} + c_2}
           {c_1'\,t_{\text{go}} + c_2'}.
  \label{eq:bts}
\end{equation}

The four coefficients are determined from four boundary conditions:
\begin{enumerate}
  \item \emph{Terminal value}: $c_2 = 0$,
        enforcing $\tan(\alpha_f + \gamma_f) = 0$ at $t_{\text{go}} = 0$.
  \item \emph{Terminal rate}: $c_1$ is set to a small constant
        (nominally $0.02$) to ensure a smooth, non-abrupt approach to
        horizontal flight.
  \item \emph{Normalisation}: $c_2' = 1$ removes the scaling ambiguity.
  \item \emph{Initial value}: $c_1'$ is computed so that the inclination
        at the current $t_{\text{go}}$ matches a blended target:
        \begin{equation}
          \gamma_{\text{blended}} = 0.7\;\gamma_{\text{current}}
            + 0.3\;\gamma_{\text{targeting}},
          \label{eq:gamma_blend}
        \end{equation}
        where $\gamma_{\text{targeting}}$ is a geometric estimate of the
        flight-path angle required to reach the target altitude over the
        remaining downrange distance.
\end{enumerate}

The bilinear form controls both the value and the rate of change of the
inclination angle at the terminal point, providing smoother orbit insertion
than the purely linear variant.

\subsubsection{Apollo Polynomial Guidance}
\label{sssec:apollo}

The most sophisticated guidance mode follows the explicit polynomial approach
originating from the Apollo programme.  Two independent channels command the
total acceleration in the downrange ($x$) and altitude ($y$) directions:
\begin{align}
  a_x(t) &= k_1\,(t - t_{\text{epoch}}) + k_2, \label{eq:apollo_ax} \\
  a_y(t) &= k_3\,(t - t_{\text{epoch}}) + k_4, \label{eq:apollo_ay}
\end{align}
where $t_{\text{epoch}}$ is the time at which the coefficients were last
recomputed.

\paragraph{Terminal constraints.}
The four coefficients $k_1,\dots,k_4$ are chosen to satisfy position and
velocity constraints at the target orbital altitude:
\begin{itemize}
  \item reach the target geocentric radius $r_f$,
  \item attain the local circular-orbit velocity
        $v_f = \sqrt{\mu / r_f}$ in the horizontal direction,
  \item achieve zero radial velocity ($\gamma_f = 0$).
\end{itemize}

The resulting expressions, derived from polynomial interpolation over the
interval $t_{\text{go}}$, are (shown for one generic channel with initial
state subscript~$i$ and final subscript~$f$):
\begin{align}
  k_1 &= \frac{6\,(v_f + v_i)\,t_{\text{go}}
          - 12\,(x_f - x_i)}{t_{\text{go}}^3},
        \label{eq:apollo_k1} \\[4pt]
  k_2 &= \frac{-2\,(v_f + 2\,v_i)\,t_{\text{go}}
          + 6\,(x_f - x_i)}{t_{\text{go}}^2}.
        \label{eq:apollo_k2}
\end{align}

Since the downrange target position is not known a~priori, it is estimated from
the current orbital elements by predicting where the unpowered trajectory would
intersect the target altitude.

\paragraph{Coefficient freezing.}
As $t_{\text{go}} \to 0$ the denominators $t_{\text{go}}^3$ and
$t_{\text{go}}^2$ drive the coefficients toward very large values.  To prevent
numerical instability the coefficients are frozen at their last computed values
once $t_{\text{go}}$ falls below a configurable threshold (nominally $10$\,s).

\paragraph{Steering-angle extraction.}
The total acceleration commanded by Eqs.~\eqref{eq:apollo_ax}--\eqref{eq:apollo_ay}
includes the gravitational contribution.  The thrust-only acceleration is
obtained by subtracting the local gravity vector, and the angle of attack is
then
\begin{equation}
  \alpha = \arctan\!\left(\frac{a_y^{\text{thrust}}}{a_x^{\text{thrust}}}\right)
           - \arctan\!\left(\frac{v_y}{v_x}\right).
  \label{eq:apollo_alpha}
\end{equation}

An optional thrust-magnitude-control mode allows the engine throttle to follow
the commanded acceleration magnitude
$|\mathbf{a}_{\text{thrust}}| = \sqrt{(a_x^{\text{thrust}})^2 +
(a_y^{\text{thrust}})^2}$, converting it to a thrust force via $F_T = m\,
|\mathbf{a}_{\text{thrust}}|$, limited by the maximum available thrust.

% -----------------------------------------------------------------------------
\subsection{Trajectory Optimisation}
\label{ssec:optimisation}

\subsubsection{Problem Formulation}
\label{sssec:opt_formulation}

The optimisation seeks the initial kick angle $\alpha_k$ that minimises the
total propellant consumption required to reach the target circular orbit:
\begin{equation}
  \min_{\alpha_k}\; J(\alpha_k)
    = m_{\text{prop,ascent}}(\alpha_k)
    + m_{\text{prop,circ}}(\alpha_k).
  \label{eq:obj}
\end{equation}

The single design variable is bounded:
$\alpha_k \in [\alpha_{\min},\,\alpha_{\max}]$ (nominally $[-5°,\,-2°]$).

Three soft constraints are enforced through a large penalty value
($J = 10^9$\,kg) applied when any condition is violated:
\begin{enumerate}
  \item propellant required for circularisation exceeds the available supply,
  \item the circularisation burn duration exceeds a maximum limit
        (nominally $15$\,s),
  \item the achieved apogee deviates from the target by more than $0.2\,\%$.
\end{enumerate}

\subsubsection{Optimisation Algorithm}
\label{sssec:opt_algorithm}

Because the objective function is non-smooth --- discrete events such as
staging, engine ignitions, and cutoffs introduce discontinuities --- a
gradient-free approach is used.  A brute-force grid search evaluates
$J(\alpha_k)$ at $N = 1000$ uniformly spaced points across
$[\alpha_{\min},\,\alpha_{\max}]$ and selects the global minimum.  No local
refinement step is applied.

The grid resolution is approximately $0.003°$, which is sufficient for
engineering accuracy.  With each trajectory evaluation taking on the order of
one second, the full optimisation completes in roughly $15$--$20$~minutes.

A key architectural feature is that the optimisation loop is entirely
independent of the guidance mode.  The outer loop supplies a candidate kick
angle and receives back a scalar cost; the inner simulation selects and
executes whichever guidance law is active.  Consequently the same search
procedure applies unchanged to all five guidance modes.

\subsubsection{Event-Based Engine Cutoff}
\label{sssec:seco_event}

The timing of SECO is not a free optimisation variable.  Instead, it is
determined implicitly during trajectory integration through an event function:
\begin{equation}
  g(t,\mathbf{x}) = r_{\text{apo}}(t,\mathbf{x}) - r_{\text{target}},
  \label{eq:seco_event}
\end{equation}
where $r_{\text{apo}}$ is the instantaneous apogee computed from the current
orbital elements (Eq.~\ref{eq:apsides}).  The integrator monitors $g$ during
the Stage~2 burn (above $65$\,km altitude); a sign change triggers root-finding
to locate the exact cutoff time.

The physical interpretation is direct: the engine burns until the trajectory,
if left unpowered from that instant onward, would coast to exactly the target
apogee altitude.  This event-driven approach automatically adapts the burn
duration to the trajectory shape produced by each guidance law and kick angle.

% -----------------------------------------------------------------------------
\subsection{Space Launcher Configuration}
\label{ssec:launcher}

The launch vehicle modelled in the simulator is representative of a
two-stage Falcon~9-class rocket.  Table~\ref{tab:launcher} summarises
the key parameters.

\begin{table}[htbp]
  \centering
  \caption{Launch-vehicle parameters used in the simulator.}
  \label{tab:launcher}
  \begin{tabular}{l c c c}
    \toprule
    \textbf{Parameter} & \textbf{Stage~1} & \textbf{Stage~2} & \textbf{Unit} \\
    \midrule
    Thrust (vacuum)       & 7\,600  & 934    & kN   \\
    Specific impulse      & 283     & 348    & s    \\
    Structural mass       & 25\,600 & 3\,900 & kg   \\
    Propellant mass       & 395\,700& 92\,670& kg   \\
    \bottomrule
  \end{tabular}
\end{table}

Stage separation occurs $3$\,s after MECO, and the second-stage engine ignites
$8$\,s after MECO, giving a $5$\,s coast interval between stages.

The aerodynamic reference area is $A = 10.52$\,m$^2$ with a constant drag
coefficient $C_D = 0.3$.

\subsubsection{Mission Parameters}

The baseline mission targets a $500$\,km circular orbit at an inclination of
$51.6°$ (representative of the International Space Station orbit), launching
from a latitude of $28.5°$ (Kennedy Space Center).  The resulting initial mass
at lift-off is approximately $517.9$\,t (sum of both stages, with zero payload
in the reference configuration).  Table~\ref{tab:mission} lists the nominal
mission parameters.

\begin{table}[htbp]
  \centering
  \caption{Baseline mission parameters.}
  \label{tab:mission}
  \begin{tabular}{l c c}
    \toprule
    \textbf{Parameter} & \textbf{Value} & \textbf{Unit} \\
    \midrule
    Target orbital altitude   & 500    & km   \\
    Target inclination        & 51.6   & deg  \\
    Launch latitude           & 28.5   & deg  \\
    Kick-angle search range   & $[-5,\,-2]$ & deg \\
    Guidance update rate      & 0.5    & s    \\
    Integration tolerance     & $10^{-8}$ & --- \\
    Maximum integration step  & 1.0    & s    \\
    Output time resolution    & 0.01   & s    \\
    \bottomrule
  \end{tabular}
\end{table}
